\section{Demostración de \texorpdfstring{\(2^{\aleph_0}=\aleph_1\)}{2 elevado a aleph cero igual a aleph uno} en el continuo HMT}
\label{sec:cierre-cardinal-nativo}

Esta sección reúne en una demostración contigua la prolongación, la doble
proyección y la profundidad arbitraria; construye después la clausura
semántica, el levantamiento uniforme, la maquinaria genealógica y el cierre
cardinal. La jerarquía conjuntista se identifica únicamente al término de la
construcción y no puede ser leída como origen causal. La dependencia causal es:

\[
 \mathrm{APP}\longrightarrow\mathrm{TRIT}\longrightarrow\mathrm{TPK}
 \longrightarrow X_\infty^{\rm enr}
 \longrightarrow\mathcal C_{\rm cont}^{\rm disc}
 \longrightarrow m_\infty
 \longrightarrow\mathfrak G_{\rm HMT}(h,m_\infty).
\tag{32}\label{eq:ch-cadena-causal}
\]

Primero se construyen el estado, sus prolongaciones, el continuo y el registro genealógico.
Después se prolonga \textbf{esa genealogía producida} mediante una regla semántica
explícita a todos sus objetos hereditarios. La identificación posterior con
una jerarquía conjuntista conocida no selecciona las semillas, los operadores,
el registro genealógico, la recta ni sus proyecciones. Sí permite contrastar, con total
transparencia, la fuerza lógico-cardinal exacta de la regla semántica adoptada.

\subsection{Prolongación de estados y extensión semántica transfinita}

Sea

\[
 \Sigma_{\rm HMT}
 =\bigl(\dot R_{\rm APP},\dot R_{\rm TRIT},\dot R_{\rm TPK},
 \dot R_{U},\dot R_{\Gamma_9},\dot R_{\rm Ext},\dot R_\rho,
 \dot R_{\pi},\dot R_{\varphi},\dot R_e,\dot R_{\alpha},
 \dot R_{\rm ar},\dot R_{\rm exc}\bigr)
\tag{33}\label{eq:ch-signatura}
\]

la signatura \textbf{relacional} numerable cuyo código efectivo es \(h\). Cada
\(\dot R_S\) es el grafo del operador \(S\) sobre los códigos donde está
definido; no se presupone que una función externa sea total dentro de cada
nivel. Los grafos \(\dot R_{\pi},\dot R_{\varphi},\dot R_e,\dot R_{\alpha}\)
son los lectores de salida construidos por la generación coinductiva y el
cierre dodecafásico; no reciben como argumentos los valores convencionales de
las constantes. Su presencia en \(\Sigma_{\rm HMT}\) asegura que la clausura
conserve el estado completo: integra simultáneamente las rutas que expresan
valor y la incidencia que se prolonga hacia la realización excepcional.

En cada horizonte, el código del estado registra las coordenadas APP--TRIT--TPK,
la ruta, la hoja, el acarreo, la frontera y la memoria; en sus etapas de
representación registra asimismo \(P_N,\Phi_N,E_N,A_N\), donde
\(A_N\) es el prefijo de \(\alpha_{\rm HMT}\) producido por el mismo estado.
Estas coordenadas no seleccionan la rama desde fuera: son parte de la historia
que se codifica. Fijada una numeración primitivo-recursiva de las coordenadas finitas de
un estado, toda historia inversa
\(x_\infty=(x_N)_{N\geq1}\in X_\infty^{\rm enr}\) determina el código
\[
 \operatorname{code}(x_\infty)
 :=\{\langle N,j,c\rangle:(x_N)_j\text{ tiene código }c\}\subseteq\omega.
\]
Las ecuaciones de coherencia
\(\rho_{N+1,N}(x_{N+1})=x_N\) permiten decodificar de manera única la
historia a partir de ese real. Escribimos \(m_\infty\) para el código de la
historia realizada y \(m_\infty(n)\) para su función característica bajo la
numeración fijada. El código operatorio \(h\) y este registro completo se
reúnen, mediante una función primitiva recursiva de emparejamiento, en el real

\[
 u_{h,m}:=
 \{\langle0,\langle n,h(n)\rangle\rangle:n<\omega\}
 \cup
 \{\langle1,\langle n,m_\infty(n)\rangle\rangle:n<\omega\}.
\tag{33a}\label{eq:ch-codigo-conjunto}
\]

Las dos proyecciones primitivas recuperan \(h\) y \(m_\infty\) desde
\(u_{h,m}\).

La primera extensión es la prolongación métrica y operatoria de los estados
finitos. Para \(N\ge1\), sea
\(\widehat{\mathcal X}^{\rm enr,raw}_{6N}\) el conjunto de estados
enriquecidos brutos de longitud \(6N\), y defínase conjuntistamente el
subsistema superviviente por
\[
\begin{aligned}
 X_N
 &:=\widehat{\mathcal X}^{\rm enr,surv}_{6N}\\
 &:=\Bigl\{x\in\widehat{\mathcal X}^{\rm enr,raw}_{6N}:\\[-.2ex]
 &\qquad\forall M>N\ \exists z\in
   \widehat{\mathcal X}^{\rm enr,raw}_{6M}\;\text{ tal que}\\[-.2ex]
 &\qquad\operatorname{tr}_{6M,6N}(z)=x\ \wedge\
   z\text{ satisface toda transición TPK intermedia}\Bigr\}.
\end{aligned}
\]
y póngase
\[
 \rho_{N+1,N}:=\operatorname{tr}_{6N+6,6N}\!\upharpoonright X_{N+1},
\]
y
\[
 \operatorname{Ext}_N(x):=
 \{y\in X_{N+1}:(x,y)\in
 \operatorname{Ext}_{6N+6,6N}\}.
\]
El nivel raíz se incluye sin excepción implícita:
\[
 X_0:=\{\ast\},\qquad
 \rho_{1,0}:X_1\to X_0\text{ es constante},\qquad
 \operatorname{Ext}_0(\ast):=X_1.
\]
El superíndice \(\mathrm{surv}\) denota el
árbol podado de los estados que poseen descendientes TPK compatibles a todo
horizonte finito. La ley de prolongación TPK establecida en las secciones
precedentes prueba que las semillas
emitidas pertenecen a ese árbol; la ramificación finita y el lema de Kőnig
convierten la compatibilidad a todo horizonte en una rama infinita. En
particular,

\[
 \forall N\ge0\;\forall x\in X_N\quad
 \varnothing\ne\operatorname{Ext}_N(x)
 \subseteq\rho_{N+1,N}^{-1}(x).
\tag{34}\label{eq:ch-extension-no-vacia}
\]

Por el mismo lema de árbol finitamente ramificado, toda semilla superviviente
posee una prolongación compatible

\[
 x_\infty=(x_N)_{N\ge1}\in
 X_\infty^{\rm enr}:=\varprojlim_N(X_N,\rho_{N+1,N}).
\tag{35}\label{eq:ch-limite-inverso}
\]

La segunda extensión actúa \textbf{después} sobre el código completo de esa
historia realizada.
No se infiere de la mera existencia del operador finito
\(\operatorname{Ext}_N\): formaliza a escala conjuntista el principio HMT de
que sólo son objetos del dominio íntegro los que poseen una derivación
genealógica desde el estado y sus reglas. La extensión semántica transfinita
es:

\[
 \begin{aligned}
  \mathfrak G_0(h,m_\infty)
  &:=\operatorname{TC}\bigl(\{\omega,u_{h,m},h,m_\infty\}\bigr),\\
  \mathfrak G_{\beta+1}(h,m_\infty)
  &:=\operatorname{Def}_{\Sigma_{\rm HMT}}
       \bigl(\mathfrak G_\beta(h,m_\infty)\bigr),\\
  \mathfrak G_\lambda(h,m_\infty)
  &:=\bigcup_{\beta<\lambda}\mathfrak G_\beta(h,m_\infty)
       \quad(\lambda\text{ límite}),\\
  \mathfrak G_{\rm HMT}(h,m_\infty)
  &:=\bigcup_{\beta\in\operatorname{Ord}}
       \mathfrak G_\beta(h,m_\infty).
 \end{aligned}
\tag{36}\label{eq:ch-clausura-transfinita}
\]

La hipótesis del continuo y una biyección objetivo no se introducen entre las
condiciones que definen el operador sucesor; el operador está dado
prospectivamente por

\[
 \operatorname{Def}_{\Sigma_{\rm HMT}}(M)
 :=M\cup
 \Bigl\{
   \{x\in M:\langle M,\in,
       (\dot R_S\cap M^{k_S})_{S\in\Sigma_{\rm HMT}}\rangle
       \models\varphi(x,\bar p)\}:
   \ulcorner\varphi\urcorner\in\omega,
   \ \bar p\in M^{<\omega}
 \Bigr\}.
\tag{37}\label{eq:ch-operador-def}
\]

Es decir: expresa simultáneamente todos los subconjuntos definibles por una
fórmula finita del lenguaje relacional HMT y parámetros que la genealogía ya
produjo. No
consulta un valor real objetivo, una constante convencional, un cardinal
objetivo ni un modelo exterior. La fórmula~\eqref{eq:ch-operador-def} define el cierre transfinito
sobre las semillas, los operadores, el estado enriquecido, el registro y la
doble proyección construidos previamente.

\begin{lemma}[Codificación conjunta del programa HMT y del registro realizado]
\label{thm:ch-6}
El par formado por
el código operatorio \(h\) y el registro \(m_\infty\) se codifica por un único
real \(u_{h,m}\subseteq\omega\), y toda fórmula de \(\Sigma_{\rm HMT}\) se
traduce uniformemente a una fórmula del lenguaje \(\{\in,u_{h,m}\}\).

\end{lemma}
\begin{proof}
Tanto \(h\) como \(m_\infty\) son sucesiones de naturales: \(h\)
codifica la signatura, las tablas finitas y los programas de los operadores;
\(m_\infty\) codifica una sola historia realizada, no el espacio completo de
historias. La fórmula~\eqref{eq:ch-codigo-conjunto} y sus dos decodificadores prueban la primera
afirmación. Cada símbolo de~\eqref{eq:ch-signatura} designa su \textbf{relación de grafo codificada en \(h\)}, nunca el grafo externo de todas sus posibles evaluaciones. APP,
TRIT, \(U_t\), \(\Gamma_9\), las restricciones \(\rho\) y las extensiones
finitas tienen grafos primitivo-recursivos sobre códigos enteros y, por tanto,
fórmulas \(\Delta_0\) absolutas entre niveles transitivos. Las dos
representaciones se expresan como relaciones de grafo mediante las fórmulas

\[
 \begin{aligned}
 \dot R_{\rm ar}(x,y)
 &\;\Longleftrightarrow\;
 \forall k\in\omega\;\exists N\;\forall n\ge N\quad
 |\mu_n(x)-y|<2^{-k},\\
 \dot R_{\rm exc}(x,c,z)
 &\;\Longleftrightarrow\;
 \forall N\in\omega\quad z\!\upharpoonright N
   =E_N(x\!\upharpoonright N,c\!\upharpoonright N),
 \end{aligned}
\tag{37c$_0$}\label{eq:ch-grafos-publicacion}
\]

donde \(\mu_n\) y \(E_N\) son las funciones racionales/finitas cuyos códigos
figuran en \(h\). Se sustituye inductivamente cada fórmula atómica que contiene
un símbolo HMT por su fórmula de grafo con parámetro \(u_{h,m}\). Los
conectivos y cuantificadores se conservan. Esta inducción sobre la complejidad
sintáctica produce una traducción \(\varphi\mapsto\varphi^*\) tal que, para
todo nivel transitivo \(M\) de~\eqref{eq:ch-clausura-transfinita} que contiene los parámetros,

\[
 \langle M,\in,(\dot R_S\cap M^{k_S})_S\rangle
   \models\varphi(\bar a)
 \quad\Longleftrightarrow\quad
 M\models\varphi^*(\bar a;u_{h,m}).
\tag{37c}\label{eq:ch-eliminacion-uniforme}
\]

Así, \eqref{eq:ch-operador-def} no puede consultar un oráculo, una tabla decimal objetivo ni una
relación no codificada que introduzca incontables objetos en un solo paso.
\end{proof}

La semántica HMT se define ahora \textbf{independientemente de identificarla con \(\mathfrak G\)}. Sea \(\mathcal T_0^{\rm HMT}\) el conjunto de nombres
canónicos de los elementos de
\(\mathfrak G_0(h,m_\infty)=
\operatorname{TC}(\{\omega,u_{h,m},h,m_\infty\})\). En particular, el nivel
inicial contiene el registro realizado \(m_\infty\), no todas las historias de
\(X_\infty^{\rm enr}\) como parámetros primitivos.
Los términos admisibles son los árboles bien fundados obtenidos por

\[
 \begin{aligned}
 \mathcal T^{\rm HMT}_{\beta+1}
 &:={\mathcal T}^{\rm HMT}_\beta\cup
 \{\langle\beta,\ulcorner\varphi\urcorner,
       t_0,\ldots,t_{k-1}\rangle:
       t_i\in\mathcal T^{\rm HMT}_\beta\},\\
 \mathcal T^{\rm HMT}_\lambda&:=
       \bigcup_{\beta<\lambda}\mathcal T^{\rm HMT}_\beta,
 \qquad
 \mathcal T^{\rm HMT}:=
       \bigcup_{\beta\in\operatorname{Ord}}\mathcal T^{\rm HMT}_\beta.
 \end{aligned}
\tag{37a}\label{eq:ch-terminos}
\]

El valor de un término sucesor es el subconjunto del nivel anterior definido
por \(\varphi\) en la estructura relacional inducida, con los valores de
\(t_0,\ldots,t_{k-1}\) como parámetros. Se define, antes de cualquier
comparación conjuntista,

\[
 \operatorname{Ob}_{\rm sem}^{\rm HMT}(h,m_\infty)
 :=\{\llbracket t\rrbracket:t\in\mathcal T^{\rm HMT}\}.
\tag{37b}\label{eq:ch-objetos-semanticos}
\]

Ésta es la regla de \textbf{exhaustividad generativa} adoptada por HMT: todo
objeto del dominio íntegro debe poseer uno de esos árboles de derivación y todo
árbol admisible expresa un objeto. No es consecuencia de la ramificación
finita. La hipótesis del continuo no se introduce en esta regla; se obtiene
mediante la demostración desarrollada en el artículo. Tampoco se introducen
como datos una enumeración de los reales ni una biyección objetivo.

\begin{theorem}[Equivalencia entre la semántica de términos y la jerarquía genealógica]
\label{thm:ch-6a}
La evaluación de
los términos anteriores satisface

\[
 \operatorname{Ob}_{\rm sem}^{\rm HMT}(h,m_\infty)
 =\mathfrak G_{\rm HMT}(h,m_\infty).
\tag{37d}\label{eq:ch-equivalencia-semantica}
\]

\end{theorem}
\begin{proof}
Por inducción transfinita. En el nivel inicial las dos clases
coinciden. Si todo valor de un término de
\(\mathcal T^{\rm HMT}_\beta\) pertenece a
\(\mathfrak G_\beta\), la semántica del término sucesor es uno de los
subconjuntos de~\eqref{eq:ch-operador-def}. Recíprocamente, cada subconjunto de~\eqref{eq:ch-operador-def} viene dado por
un código de fórmula y una tupla finita de parámetros; por la hipótesis
inductiva esos parámetros poseen términos y~\eqref{eq:ch-terminos} forma el término que lo
denota. En un límite ambas construcciones toman la misma unión.
\end{proof}

\Needspace{9\baselineskip}
El codificador de estados es compatible con los truncamientos y envía cada
prolongación TPK a un término sucesor. Sea
\(\mathcal T_N^{\rm st}\subset\mathcal T^{\rm HMT}\) la subclase de términos
que codifica estados completos de profundidad \(N\), y sea
\(\tau_{N+1,N}:\mathcal T_{N+1}^{\rm st}\to\mathcal T_N^{\rm st}\) el borrado
del último bloque junto con el truncamiento correspondiente del registro genealógico,
acarreo y memoria. Si \(J_N:X_N\to\mathcal T_N^{\rm st}\) es el codificador
completo, la compatibilidad de \(\operatorname{Ext}_N\) prueba la identidad
correctamente tipada

\[
 \forall x\in X_N\ \forall y\in\operatorname{Ext}_N(x)\qquad
 \tau_{N+1,N}\bigl(J_{N+1}(y)\bigr)=J_N(x).
\tag{37e$_0$}\label{eq:ch-naturalidad-codificador}
\]

Defínase dentro de la semántica
\[
 \operatorname{Succ}^{\rm TPK}_{N+1}(J_N(x)):=
 \{t\in\mathcal T_{N+1}^{\rm st}:\exists y\in X_{N+1}\,
   (y\in\operatorname{Ext}_N(x)\ \wedge\ t=J_{N+1}(y))\}.
\]
Elíjase \(\beta\) tal que \(\mathfrak G_\beta\) contenga \(x\), \(X_N\),
\(X_{N+1}\), los grafos finitos de
\(\operatorname{Ext}_N,J_N,J_{N+1}\) y sus códigos en \(h\). Entonces la
fórmula primitivo-recursiva del grafo de \(\operatorname{Ext}_N\) da

\[
 J_{N+1}[\operatorname{Ext}_N(x)]
 =\operatorname{Succ}^{\rm TPK}_{N+1}(J_N(x)),
 \qquad
 J_{N+1}[\operatorname{Ext}_N(x)]
 \in\mathfrak G_{\beta+1},
\tag{37e}\label{eq:ch-sucesores-semanticos}
\]

y ambos miembros son formados por la fórmula de grafo de la misma extensión
TPK. Las ecuaciones~\eqref{eq:ch-naturalidad-codificador}--\eqref{eq:ch-sucesores-semanticos} son el enlace semántico explícito: la
subclase de términos sucesores es exactamente la imagen de las prolongaciones
admisibles y su truncamiento conmuta. No se identifica la fibra completa de
\(\rho\) con \(\operatorname{Ext}_N\), ni se identifica
\(\operatorname{Def}\) con TPK por nomenclatura.

\begin{lemma}[Transitividad, axiomas y buen orden del modelo genealógico]
\label{thm:ch-6b}
\(\mathfrak G_{\rm HMT}(h,m_\infty)\) es una clase transitiva, contiene todos
los ordinales, satisface la teoría de conjuntos de Zermelo--Fraenkel (ZF) y
posee un buen orden global definible; satisface,
por tanto, elección y todos los axiomas usados en la comparación cardinal posterior.

\end{lemma}
\begin{proof}
Abreviemos \(B=\mathfrak G_0\) y \(u=u_{h,m}\). La inducción
sobre~\eqref{eq:ch-clausura-transfinita} prueba simultáneamente que cada nivel es transitivo y que
\(\mathfrak G_\alpha\in\mathfrak G_{\alpha+1}\); si todos los ordinales
menores que \(\alpha\) están presentes, su conjunto es definible en el nivel
correspondiente y expresa \(\alpha\). Por ello la unión contiene todos los
ordinales. Extensionalidad y Fundación se heredan del universo ambiente;
Vacío e Infinito están en \(B\). Elegidos un nivel que contenga los parámetros,
las fórmulas que definen \(\{a,b\}\), \(\bigcup a\) y los subconjuntos por
Separación los expresan en un sucesor.

Para Colección--Reemplazo, sea \(F\) una relación funcional definible en
\(\mathfrak G\) sobre el conjunto \(a\). El esquema de reflexión relativo a
la jerarquía definible \((\mathfrak G_\xi,\in,u_{h,m})_{\xi\in\mathrm{Ord}}\),
aplicado en la metateoría ambiente a la lista finita de fórmulas que define
\(F\), produce un nivel
\(\mathfrak G_\theta\) que contiene \(a\), los parámetros y es correcto para
esas fórmulas. Los rangos de los testigos de \(F[a]\) están entonces acotados
por \(\theta\); Separación sobre \(\mathfrak G_\theta\) forma la imagen en
\(\mathfrak G_{\theta+1}\). Para Potencia no se presupone el conjunto interno
que se quiere construir. Fijado \(a\in\mathfrak G_\gamma\), se considera en
la metateoría ambiente el conjunto ya existente \(\mathcal P^V(a)\). Para
cada \(b\in\mathcal P^V(a)\cap\mathfrak G\), sea \(r(b)\) su primera etapa
genealógica. Como \(\mathfrak G\) es una clase definible, Separación en la
metateoría ambiente forma primero el conjunto
\(\mathcal P^V(a)\cap\mathfrak G\). Reemplazo \textbf{en la metateoría ambiente}
aplicado a ese conjunto acota los ordinales \(r(b)\) por algún \(\theta\).
En consecuencia,
\[
 \mathcal P^{\mathfrak G}(a)
 =\{b\in\mathfrak G_\theta:b\subseteq a\},
\]
y el miembro derecho se expresa por Separación en
\(\mathfrak G_{\theta+1}\). No introduce en \(\mathfrak G\) los subconjuntos
exteriores a la semántica~\eqref{eq:ch-objetos-semanticos}. Éste es el argumento no circular de
acotación por rangos.

Finalmente, a cada objeto se le asigna su primer nivel de nacimiento y su
menor \textbf{nombre de derivación}: un árbol de derivación bien fundado y
finitamente ramificado (y, por~\eqref{eq:ch-terminos}, finito para cada nombre individual), etiquetado por un
ordinal de nivel, un código de fórmula y la tupla de nombres de sus parámetros.
Por recursión transfinita se ordenan esos nombres: primero la etiqueta ordinal,
después el código de fórmula y después, lexicográficamente, los nombres de los
parámetros ya ordenados. El menor nombre de cada valor define una relación de
clase \(<_{\mathfrak G}\) que compara todos los objetos. Es un buen orden
global definible y, tomando el menor elemento de cada conjunto no vacío,
prueba Elección.
\end{proof}

\paragraph{Reflexión por clausura de testigos y construcción de la imagen.}
El paso de Colección--Reemplazo del lema~\ref{thm:ch-6b} admite la siguiente
demostración explícita. Se mantiene la jerarquía ya generada a partir del
código operatorio y del registro realizado; abreviamos
\(\mathfrak G=\mathfrak G_{\rm HMT}(h,m_\infty)\).
La argumentación utiliza Separación y Reemplazo en la metateoría ambiente,
sin presuponer estos esquemas en \(\mathfrak G\).

Sea \(\Phi\) una colección finita de fórmulas. Se traducen primero sus
cuantificadores universales mediante negación y cuantificación existencial,
y se cierra la colección resultante bajo subfórmulas. Para cualquier nivel
inicial existe un ordinal límite posterior \(\theta\) tal que
\[
 \mathfrak G_\theta\models\varphi(\bar a)
 \quad\Longleftrightarrow\quad
 \mathfrak G\models\varphi(\bar a)
 \qquad
 (\varphi\in\Phi,\ \bar a\in\mathfrak G_\theta^{<\omega}).
\]
La notación de satisfacción de la clase se interpreta, para cada fórmula
fijada, relativizando sus cuantificadores a la clase definible
\(\mathfrak G\); no se introduce un predicado uniforme de verdad.

\begin{proof}
Fijemos \(\alpha\). Para cada subfórmula existencial
\(\exists y\,\psi(y,\bar x)\) de \(\Phi\), las tuplas
\(\bar a\in\mathfrak G_\alpha^{<\omega}\) de la aridad correspondiente
que poseen un testigo en \(\mathfrak G\) forman un conjunto por Separación
metateórica. A cada una se asigna el menor ordinal \(\beta\) para el que
algún \(y\in\mathfrak G_\beta\) satisface
\(\mathfrak G\models\psi(y,\bar a)\). El ordinal existe porque
\(\mathfrak G\) es la unión de sus niveles. Reemplazo en la metateoría
acota esas primeras etapas de testigos. La finitud de \(\Phi\) permite
elegir definiblemente una sola cota \(f(\alpha)>\alpha\), válida para
todas sus subfórmulas existenciales. Por tanto, cada tupla considerada tiene
\emph{al menos un testigo} en \(\mathfrak G_{f(\alpha)}\).

Tomamos \(\alpha_0\) suficientemente grande para contener los parámetros
iniciales y ponemos
\[
 \alpha_{n+1}=f(\alpha_n),\qquad
 \theta=\sup_{n<\omega}\alpha_n.
\]
Por continuidad de la jerarquía, toda tupla finita de
\(\mathfrak G_\theta\) pertenece a algún \(\mathfrak G_{\alpha_n}\).
Si una fórmula existencial de \(\Phi\) con esa tupla es verdadera en
\(\mathfrak G\), posee un testigo en
\(\mathfrak G_{\alpha_{n+1}}\subseteq\mathfrak G_\theta\).
La inducción sobre las fórmulas prueba ahora la equivalencia: las atómicas
se conservan en la estructura inducida; los conectivos preservan la
equivalencia; el paso existencial utiliza ese testigo en una dirección y el
mismo testigo, visto en \(\mathfrak G\), en la otra. La traducción previa
de los universales incluye también las fórmulas existenciales que expresan
sus posibles contraejemplos. Queda probada la reflexión finita requerida.
\end{proof}

Sea ahora \(F(x,y,\bar p)\) funcional sobre \(a\) en \(\mathfrak G\).
Aplicamos la construcción a las fórmulas que expresan \(F\), la existencia
de sus testigos y su imagen, con \(a,\bar p\in\mathfrak G_\theta\).
La transitividad da \(a\subseteq\mathfrak G_\theta\). Cada
\(x\in a\) posee su valor dentro de ese nivel, y la reflexión identifica
la imagen exacta con
\[
 b=\{y\in\mathfrak G_\theta:
       \mathfrak G_\theta\models
       \exists x\in a\ F(x,y,\bar p)\}.
\]
Este conjunto es definible sobre \(\mathfrak G_\theta\), con parámetros
del mismo nivel. Por la operación sucesora de la jerarquía,
\(b\in\operatorname{Def}(\mathfrak G_\theta)
\subseteq\mathfrak G_{\theta+1}\). Así se obtiene Reemplazo en
\(\mathfrak G\). Si se conserva la existencia de testigos y se omite la
unicidad, la misma construcción proporciona un conjunto de testigos para
Colección.

\paragraph{Aridad de los nombres y selección interna.}
En el buen orden del lema~\ref{thm:ch-6b}, la enumeración de la base
asigna a cada elemento su menor código natural. En cada sucesor se ordenan
primero los códigos de fórmula; fijado uno, queda fijada la aridad de su
tupla de parámetros. El orden lexicográfico se aplica entonces a un
\emph{producto finito de aridad fija} de los buenos órdenes anteriores.
El menor nombre representa cada valor nuevo y los niveles límite conservan
el orden por primera aparición. De este modo queda especificado el orden
de las tuplas, sin mezclar aridades distintas en un único orden
lexicográfico. La selección del menor elemento de cada conjunto no vacío,
junto con Reemplazo ya obtenido, forma internamente las funciones de
elección utilizadas después. Esta construcción de nombres precede a las
inyecciones cardinales y no recibe ninguna biyección del continuo como
dato.


\subsection{Elevación TPK de la fibra ternaria y partición cilíndrica completa}
\label{subsec:elevacion-catalogal-ch}

La fibra de \(729=3^6\) subcilindros de refinamiento no procede del cociente de transferencia
\(\mathcal K_{\rm ph}\). Ese operador es una representación posterior de memoria
reducida y no puede generar historias. La cobertura se construye directamente
en el sistema \((H_k,\rho_k,\operatorname{Ext}_k)\) mediante las tres
prolongaciones de una vuelta de la holonomía nonádica.

\subsubsection{Inversión de Hensel en la torre aritmética y realización por historias}
\label{sec:ch-capacidad-hensel}

Antes de levantar las ramas a historias enriquecidas, conviene recuperar el
argumento finito que determina la capacidad de la torre aritmética. Se
incorpora la prueba completa del transductor no filtrado de la matriz; su
dominio se distingue del árbol superviviente.

\begin{proposition}[Conjugación de la torre de Hensel no filtrada]
\label{prop:ch-hensel-bruto}
Sean \(p\) primo, \(d\geq1\) y \(A\in\operatorname{GL}_d(\mathbb Z)\).
Para vectores fila, tomando representantes de residuos en
\(\{0,\ldots,p-1\}^d\), se define
\[
 x_tA=u_t+px_{t+1},\qquad u_t=x_tA\bmod p.
\]
La aplicación
\[
 \Psi_T:(\mathbb Z/p^T\mathbb Z)^d\longrightarrow
             (\mathbb F_p^d)^T,\qquad
 x_0\longmapsto(u_0,\ldots,u_{T-1})
\]
es biyectiva, con inversa
\[
 x_0\equiv\sum_{t=0}^{T-1}p^tu_tA^{-(t+1)}\pmod{p^T}.
\]
Las biyecciones conmutan con las reducciones y dan un homeomorfismo
\(\Psi:\mathbb Z_p^d\to(\mathbb F_p^d)^{\mathbb N}\), donde
\(\mathbb Z_p:=\varprojlim_T\mathbb Z/p^T\mathbb Z\). La actualización
\(x_t\mapsto x_{t+1}\) corresponde a suprimir el primer bloque.
\end{proposition}

\begin{proof}
La inversa entera de \(A\) existe porque su determinante es \(\pm1\).
Iterando \(x_t=(u_t+px_{t+1})A^{-1}\), se obtiene
\[
 x_0=\sum_{t=0}^{T-1}p^tu_tA^{-(t+1)}+p^Tx_TA^{-T}.
\]
Módulo \(p^T\), una palabra determina una única clase inicial. Dada una
palabra arbitraria, fijar \(x_T=0\) y reconstruir hacia atrás produce una
clase que la emite, lo que prueba la sobreyectividad. La misma fórmula
prueba compatibilidad al truncar. El paso al límite inverso da el
homeomorfismo y la relación con el desplazamiento de la cinta.
\end{proof}

Para \(p=3\) y \(d=6\), el alfabeto de bloques tiene \(3^6=729\)
elementos. Esta capacidad aritmética no se identifica por definición con
la realizabilidad de cada bloque tras cada prefijo superviviente. El
enlace con las aristas \(\operatorname{Ext}\) es precisamente el cometido
del lema~\ref{lem:elevacion-catalogal-nonadica}. Se mantienen así las dos
operaciones que distingue la matriz: representación en la torre bruta y
realización bajo las compatibilidades APP--TRIT--TPK.


\subsubsection{Realización arista por arista de las tres prolongaciones nonádicas}
\label{subsubsec:tres-vueltas-ext-ch}

Esta subsección especifica el modelo autónomo generado por la relación de
extensión TPK que interviene en la prueba cardinal. No se infiere la existencia de
una arista por coincidencia de tipos: se construyen sus estados inicial y final,
su relación APP, su transporte de fibra y la actualización unaria del registro.
La construcción utiliza exclusivamente las dos hojas de APP, el selector TRIT
y la composición \(\mathsf{Sel}\)--\(\mathsf{Tra}\)--\(\operatorname{Upd}\) del TPK.

\paragraph{Células completas y tres reducciones balanceadas.}
En la coordenatización \(D_9=\{1,\ldots,9\}\), mantendremos siempre la célula APP
completa
\[
 a_{p,q}:=(p,q;R^\Sigma_{p,q},q^\Sigma_{p,q},
                 R^\Pi_{p,q},q^\Pi_{p,q}),
\]
con
\[
 p+q=R^\Sigma_{p,q}+9q^\Sigma_{p,q},\qquad
 pq=R^\Pi_{p,q}+9q^\Pi_{p,q}.
\]
La hoja activa se registra en una coordenada separada; no se escinde
\(a_{p,q}\) en una supuesta célula aditiva y otra multiplicativa. Para
\(r\in D_9\), sean
\begin{equation}
 a(r):=
 \begin{cases}
  a_{1,9},&r=1,\\
  a_{1,r-1},&2\le r\le9,
 \end{cases}
 \qquad
 \tau(r):=M_{\rm ph}(r),
 \qquad
 m(r):=\tau(r)\in\{-1,0,+1\}.
 \label{eq:celula-completa-fase-ch}
\end{equation}
El residuo aditivo de \(a(r)\) proporciona aquí un representante APP de la
fase; esto no convierte \(\Sigma\) en hoja activa universal. El selector
dinámico activa \(\Sigma\) en \(r\in\{1,4,7\}\), activa \(\Pi\) en
\(r\in\{2,5,8\}\) y conserva ambas hojas, sin desplazar cursor, en
\(r\in\{3,6,9\}\). En los tres casos se transportan conjuntamente
\(R^\Sigma,q^\Sigma,R^\Pi,q^\Pi\). Se tiene
\[
 (m(1),\ldots,m(9))=(1,-1,0,1,-1,0,1,-1,0).
\]
Las tres parejas ordenadas por la base nonádica son
\begin{equation}
 \mathfrak p_{+1}:=(1,4),\qquad
 \mathfrak p_{-1}:=(2,5),\qquad
 \mathfrak p_{0}:=(3,6).
 \label{eq:tres-pares-internos-ch}
\end{equation}
Estas parejas son la sección normalizada que toma las dos primeras
ocurrencias de cada tipo TRIT bajo el desplazamiento \(r\mapsto r+3\); el
tercer bloque de fases \(7,8,9\) conserva el control de cierre de la vuelta.
Sus dos miembros pertenecen a una misma fibra de \(M_{\rm ph}\). La
reducción balanceada del capítulo TRIT da, sin introducir un coeficiente
objetivo,
\begin{equation}
 \begin{array}{c|c|c|c}
 \varepsilon&m(\mathfrak p_\varepsilon^{(1)})
                  +m(\mathfrak p_\varepsilon^{(2)})
   &r_3\bigl(2\varepsilon\bigr)&
   \chi(\varepsilon,\varepsilon)\\ \hline
 +1&1+1=(-1)+3\cdot1&-1&+1\\
 0&0+0=0+3\cdot0&0&0\\
 -1&(-1)+(-1)=1+3\cdot(-1)&+1&-1
 \end{array}
 \label{eq:tres-carry-derivados-ch}
\end{equation}
Por tanto,
\begin{equation}
 \chi(\varepsilon,\varepsilon)
 =\frac{2\varepsilon-r_3(2\varepsilon)}3=\varepsilon.
 \label{eq:carry-no-parametrico-ch}
\end{equation}
El signo de la rama es así la salida del cociclo balanceado aplicado a dos
emisiones TRIT producidas por APP. Cada pareja proporciona una de las tres
salidas de acarreo. Por claridad, indexaremos después la pareja mediante
\(\varepsilon(\mathfrak p):=
\chi(m(\mathfrak p^{(1)}),m(\mathfrak p^{(2)}))\); el símbolo
\(\varepsilon\) no suministra a \(\operatorname{Upd}^{\rm cal}\) un acarreo libre.

La arquitectura de fase, el levantamiento balanceado y el transporte de
memoria se realizan aquí mediante veintisiete incidencias locales y sus
grafos de extensión. La sección normalizada fija representantes de las
fibras de \(M_{\rm ph}\); la cobertura que se demostrará utiliza la existencia
de esos representantes y permanece invariante al sustituirlos por otra
sección compatible.

\paragraph{Datos tipados de las nueve aristas.}
Sea
\(K_{\rm Carr}:=\langle\kappa_{\rm Carr}\rangle_{\rm Ab}\) el módulo de
acarreo. Las acciones están tipadas por
\[
 H_{\rm Carr}:\mathscr P_{\rm TPK}^{\rm op}
   \longrightarrow\operatorname{End}(K_{\rm Carr}),\qquad
 Y_{\rm Carr}:\mathscr P_{\rm TPK}
   \longrightarrow\operatorname{End}(K_{\rm Carr}).
\]
En este modelo generado se toman las acciones triviales
\(H_{\rm Carr}(e)=Y_{\rm Carr}(e)=\operatorname{id}_{K_{\rm Carr}}\): la
hoja activa y la ruta se conservan como coordenadas propias del registro y no
actúan sobre el sumando libre de acarreo. Sea además
\((\mathsf M_d,\jmath_{d',d})\) el sistema directo graduado de memoria. Para
no confundir dos representantes identificados por el colímite, el estado
conserva además el grado como coordenada y se usa el portador etiquetado
\[
 \widetilde{\mathsf M}:=
 \bigsqcup_{d\ge1}\bigl(\{d\}\times\mathsf M_d\bigr),
 \qquad
 \operatorname{gr}_{\mathsf M}(d,\mu)=d.
\]
La representación en la memoria límite viene dada por
\[
\begin{aligned}
 \mathsf M&:=\varinjlim_d\mathsf M_d,&
 \jmath_d&:\mathsf M_d\longrightarrow\mathsf M,\\
 u&:=\jmath_d(u_d),&
 \iota_{\mathsf M}:\mathbb Z&\longrightarrow\mathsf M,
 \qquad n\longmapsto nu .
\end{aligned}
\]
Aquí \(\jmath_d\) es el mapa canónico del colímite y los \(u_d\) forman una
colección compatible. Su componente libre permite
el lector \(\deg_u(nu)=n\).
Las proyecciones de configuración, fase, memoria, acarreo, ruta y
supervivencia se tiparán sobre el portador calibrado construido
recursivamente más abajo; no se importará para ellas ningún espacio ambiental.
Para cada arista construida \(e\), la
componente calibrada
\(\operatorname{Upd}^{\rm cal}_e:=\mathsf{Mem}^{\rm cal}_e\) denotará la
tercera aplicación de la factorización; su ley se fija campo por campo más
abajo y no se presupone como restricción de una relación ambiental.

Fijemos la profundidad inicial \(k_0=1\), una configuración completa
\(q_\star\in\mathcal Q_{\rm obs}\) de fase uno y
\[
 q_{n,r}:=F_{\rm obs}^{9n+r-1}(q_\star),\qquad
 q_{n,10}:=q_{n+1,1}.
\]
Como el calendario observable es periódico, su valor \(q_{n,r}\) no
distingue por sí solo apariciones separadas por doce vueltas. El modelo
calibrado emplea, por ello, la elevación etiquetada
\[
 \widetilde q_{n,r}:=(q_{n,r},n,r),\qquad
 \widetilde q_{n,10}:=\widetilde q_{n+1,1},\qquad
 \pi_{\rm obs}(\widetilde q_{n,r})=q_{n,r}.
\]
La dinámica elevada
\(\widetilde F_{\rm obs}(\widetilde q_{n,r})
=\widetilde q_{n,r+1}\) proyecta a \(F_{\rm obs}\); la etiqueta no modifica
la fase física, sino que hace distintos sus niveles de aparición.
Se etiqueta la copia de la célula en cada nivel mediante
\[
 a_{n,r}:=(n,r;a(r)),\qquad
 (n,r)^+:=
 \begin{cases}
  (n,r+1),&1\le r<9,\\
  (n+1,1),&r=9.
 \end{cases}
\]
La etiqueta sólo evita identificar dos apariciones de una misma célula; las
dos divisiones euclídeas y sus cuatro coordenadas son las de \(a(r)\).
Para la vuelta \(n\ge0\) y la fase \(r\in D_9\), el símbolo
\[
 e_{\varepsilon,n,r}:\widetilde q_{n,r}
 \longrightarrow\widetilde q_{n,r+1},
 \qquad r^+:=1+(r\bmod9),
\]
ocupa el nivel
\(d_{n,r}:=k_n+r-1\), donde \(k_n:=k_0+9n\). Su registro local contiene
la célula completa \(a(r)\), la hoja activa determinada por \(M_{\rm ph}(r)\),
el tipo \(\tau(r)\),
la marca de pertenencia a \(\mathfrak p_\varepsilon\) y, al alcanzar el segundo miembro
de esa pareja, la identidad certificada
\eqref{eq:tres-carry-derivados-ch}. Definimos sus operadores únicamente a
partir de ese registro:
\begin{align}
 \mathsf C_{\mathsf M}(e_{\varepsilon,n,r})
   &=\jmath_{d_{n,r}+1,d_{n,r}},&
 \kappa_{\mathsf M}(e_{\varepsilon,n,r})
   &=\mathbf1_{\{r=9\}}u_{d_{n,r}+1},
 \label{eq:memoria-arista-ch}\\
       \operatorname{Carr}(e_{\varepsilon,n,r})
   &=d_{\varepsilon,r}\kappa_{\rm Carr},&
 d_{\varepsilon,r}
   &:=\mathbf1_{\{r=\mathfrak p_\varepsilon^{(2)}\}}
       \chi\!\left(m(\mathfrak p_\varepsilon^{(1)}),
                    m(\mathfrak p_\varepsilon^{(2)})\right).
 \label{eq:carry-arista-derivado-ch}
\end{align}
Para las identidades y para toda composición admisible se define, además,
\begin{align}
 \mathsf C_{\mathsf M}(\operatorname{id}_{\widetilde q};d)
   &:=\operatorname{id}_{\mathsf M_d},&
 \kappa_{\mathsf M}(\operatorname{id}_{\widetilde q};d)&:=0,
 \label{eq:memoria-identidad-ch}\\
 \mathsf C_{\mathsf M}(e_t\cdots e_1)
   &:=\mathsf C_{\mathsf M}(e_t)\circ\cdots\circ
      \mathsf C_{\mathsf M}(e_1),&
 \kappa_{\mathsf M}(e_2e_1)
   &:=\mathsf C_{\mathsf M}(e_2)
        \bigl(\kappa_{\mathsf M}(e_1)\bigr)
      +\kappa_{\mathsf M}(e_2).
 \label{eq:cociclo-memoria-calibrado-ch}
\end{align}
La identidad lleva el grado \(d\) porque una misma configuración observable
puede reaparecer en profundidades distintas. La última igualdad, iterada, define
\(\kappa_{\mathsf M}(e_t\cdots e_1)\) para cualquier longitud y es la ley de
cociclo del sistema directo; las inclusiones satisfacen
\(\jmath_{d+2,d+1}\jmath_{d+1,d}=\jmath_{d+2,d}\).
Además se fija
\(s(e_{\varepsilon,n,r})=\top\). Las siguientes veintisiete entradas
certifican las coordenadas que distinguen las tres prolongaciones; los campos
universales del transporte se especifican inmediatamente después. Los símbolos \(I\) y \(II\)
designan, respectivamente, la primera emisión conservada y la segunda emisión
que efectúa la reducción; un guion indica incremento de acarreo nulo, pero no
ausencia de arista.
\begin{center}
\small
\setlength{\tabcolsep}{3.2pt}
\begin{tabular}{c c cc cc cc c}
\toprule
\(r\)&\(M_{\rm ph}(r)\)&\multicolumn{2}{c}{\(\mathfrak p_{+1}\)}&
\multicolumn{2}{c}{\(\mathfrak p_0\)}&
\multicolumn{2}{c}{\(\mathfrak p_{-1}\)}&
\(\Delta\operatorname{mem}\)\\
& &marca&\(d_{+1,r}\)&marca&\(d_{0,r}\)&marca&\(d_{-1,r}\)&\\
\midrule
1&\(+1\)&I&0&--&0&--&0&0\\
2&\(-1\)&--&0&--&0&I&0&0\\
3&0&--&0&I&0&--&0&0\\
4&\(+1\)&II&\(+1\)&--&0&--&0&0\\
5&\(-1\)&--&0&--&0&II&\(-1\)&0\\
6&0&--&0&II&0&--&0&0\\
7&\(+1\)&--&0&--&0&--&0&0\\
8&\(-1\)&--&0&--&0&--&0&0\\
9&0&--&0&--&0&--&0&\(u\)\\
\bottomrule
\end{tabular}
\end{center}
Así, \(d_{\varepsilon,r}\) se calcula a partir de dos salidas del selector y
es distinto de los cocientes \(q^\Sigma,q^\Pi\) de la célula. El registro
conserva también el par marcado y el residuo balanceado; en particular, la
rama nula no se confunde con una ausencia de evento.

La ley de composición del acarreo empleada por la categoría TPK es
\begin{equation}
 \operatorname{Carr}(\operatorname{id}_{\widetilde q})=0,\qquad
 \operatorname{Carr}(e_2e_1)
 =H_{\rm Carr}(e_1)\bigl(\operatorname{Carr}(e_2)\bigr)
  +Y_{\rm Carr}(e_2)\bigl(\operatorname{Carr}(e_1)\bigr).
 \label{eq:ley-carry-correcta-ch}
\end{equation}
En identidades y caminos se prolongan functorialmente:
\[
\begin{aligned}
 H_{\rm Carr}(\operatorname{id})
   &=Y_{\rm Carr}(\operatorname{id})=\operatorname{id},\\
 H_{\rm Carr}(e_t\cdots e_1)
   &=H_{\rm Carr}(e_1)\circ\cdots\circ H_{\rm Carr}(e_t),\\
 Y_{\rm Carr}(e_t\cdots e_1)
   &=Y_{\rm Carr}(e_t)\circ\cdots\circ Y_{\rm Carr}(e_1).
\end{aligned}
\]
En este modelo todas estas acciones siguen siendo la identidad, pero
su tipado explicita qué endomorfismo actúa sobre cada sumando de
\eqref{eq:ley-carry-correcta-ch}.
En las aristas calibradas, \(H_{\rm Carr}=Y_{\rm Carr}=\operatorname{id}\),
por lo que la composición suma los incrementos ya producidos. La ecuación
\eqref{eq:ley-carry-correcta-ch}, y no una identificación entre orientación y
cociente, gobierna el registro acumulado.

\paragraph{Construcción independiente de las tuplas de transición.}
Un estado de la fibra se escribirá, en el orden de sus campos pertinentes,
\[
 x=(a;q;\tau,o,h;b_{\rm ref};
       R^\Sigma,q^\Sigma,R^\Pi,q^\Pi;c;
       \operatorname{Sig};C,\partial C;p_{\APP},p_{\TRIT};
       \lambda;\operatorname{gr}_{\mathsf M};\operatorname{mem};
       s;\operatorname{Led}).
\]
La notación de actualización \(x[f\leftarrow v]\) cambia sólo el campo
indicado. Para \(h\in\{(\Sigma,\Pi),(\Pi,\Sigma)\}\), defínase
\[
 h_r(h)=
 \begin{cases}
  (\Sigma,\Pi),&r\in\{1,4,7\},\\
  (\Pi,\Sigma),&r\in\{2,5,8\},\\
 h,&r\in\{3,6,9\}.
 \end{cases}
\]
La coordenada completa de hoja y aritmética se abrevia por
\[
 \widehat h(x):=
 \bigl(h(x),R^\Sigma(x),q^\Sigma(x),R^\Pi(x),q^\Pi(x)\bigr).
\]
Sea \(\widehat{\mathcal H}^{\rm cal}_{\widetilde q_{n,r}}\) el conjunto de
las quíntuplas
\[
 \bigl(h,R^\Sigma_{a_{n,r}},q^\Sigma_{a_{n,r}},
          R^\Pi_{a_{n,r}},q^\Pi_{a_{n,r}}\bigr),
 \qquad
 h\in\{(\Sigma,\Pi),(\Pi,\Sigma)\},
\]
cuyos cuatro datos aritméticos son las dos divisiones euclídeas de la célula
APP etiquetada \(a_{n,r}\). Esta definición directa no presupone un estado
en un portador todavía no construido. Sean
\(\mathsf{Hist}^{\rm cal}_{\widetilde q}\) el conjunto de caminos finitos
admisibles del grafo de aristas etiquetadas que terminan en
\(\widetilde q\), incluida la identidad, y
\(\mathsf B^{\rm cal}_{\widetilde q}\) el conjunto de sus sucesiones
ordenadas de extremos. Definimos
\[
 \mathsf S^{\rm cal}_{\widetilde q}:=
 \widehat{\mathcal H}^{\rm cal}_{\widetilde q}
 \times\{o_{\rightarrow},o_{\circ},o_{\leftarrow}\}
 \times K_{\rm Carr}\times
 \mathsf{Hist}^{\rm cal}_{\widetilde q}\times
 \mathsf B^{\rm cal}_{\widetilde q}\times\mathsf M .
\]
La firma calibrada no se postula: es la aplicación de reunión tipada
\[
 \operatorname{Sig}^{\rm cal}_{\widetilde q}:
 \widehat{\mathcal H}^{\rm cal}_{\widetilde q}
 \times\{o_{\rightarrow},o_{\circ},o_{\leftarrow}\}
 \times K_{\rm Carr}\times
 \mathsf{Hist}^{\rm cal}_{\widetilde q}\times
 \mathsf B^{\rm cal}_{\widetilde q}\times\mathsf M
 \longrightarrow\mathsf S^{\rm cal}_{\widetilde q},
\quad
 (\widehat h,o,c,\lambda,\partial C,\mu)
 \longmapsto(\widehat h,o,c,\lambda,\partial C,\mu).
\]
Asimismo, \(\pi_{\mathcal F}\) denota la proyección que elimina únicamente
las dos primeras coordenadas \((a;q)\) de la tupla de estado; conserva
\(\tau,o,h,b_{\rm ref}\), los cuatro datos euclídeos, acarreo, firma,
cilindro, frontera, prefijos, ruta, grado de memoria, memoria, supervivencia y
registro.
La orientación correspondiente queda fijada por
\(o(+1)=o_{\rightarrow}\), \(o(0)=o_{\circ}\) y
\(o(-1)=o_{\leftarrow}\).
La fibra APP calibrada sobre \(\widetilde q_{n,r}\) contendrá la copia etiquetada
\(a_{n,r}\) ya definida. El estado inicial queda fijado sin coordenadas libres:
\[
 \widehat h_\star:=
 \bigl((\Sigma,\Pi),R^\Sigma_{1,9},q^\Sigma_{1,9},
       R^\Pi_{1,9},q^\Pi_{1,9}\bigr).
\]
\begin{equation}
\begin{aligned}
 x_\star:=\bigl(&a_{0,1};\widetilde q_{0,1};
 +1,o_{\rightarrow},(\Sigma,\Pi);\bot;
 R^\Sigma_{1,9},q^\Sigma_{1,9},R^\Pi_{1,9},q^\Pi_{1,9};0;\\
 &\operatorname{Sig}^{\rm cal}_{\widetilde q_{0,1}}
       (\widehat h_\star,o_{\rightarrow},0,
        \operatorname{id}_{\widetilde q_{0,1}},\varnothing,
        \jmath_{k_0}(0_{\mathsf M_{k_0}}));
 \varnothing,\varnothing;
 (a_{0,1}),(+1);\operatorname{id}_{\widetilde q_{0,1}};
 k_0;0_{\mathsf M_{k_0}};\top;\varnothing\bigr).
\end{aligned}
\label{eq:estado-inicial-calibrado-ch}
\end{equation}
La tupla no contiene coordenadas indeterminadas y satisface
\(s(x_\star)=\top\). Sea \(G_0:=\{x_\star\}\).
Supuesto construido \(G_n\), fijemos \(x\in G_n\),
\(\varepsilon\in\{-1,0,+1\}\) y
\(x_{\varepsilon,0}(x):=x\). Para \(r=1,\ldots,9\), se construye primero
el registro elemental
\begin{equation}
\begin{aligned}
 \ell_{\varepsilon,n,r}:=\bigl(&e_{\varepsilon,n,r},a_{n,r},a_{(n,r)^+},
 \widetilde q_{n,r},\widetilde q_{n,r+1},M_{\rm ph}(r),\\
 &h_r(h(x_{\varepsilon,r-1}(x))),o(M_{\rm ph}(r)),\mathfrak p_\varepsilon,
 d_{\varepsilon,r},\\
 &\mathsf C_{\mathsf M}(e_{\varepsilon,n,r}),
 \kappa_{\mathsf M}(e_{\varepsilon,n,r}),
 \operatorname{Carr}(e_{\varepsilon,n,r})\bigr),
\end{aligned}
 \label{eq:registro-elemental-ext-ch}
\end{equation}
y se pone
\(\operatorname{Led}^{\rm new}_{\varepsilon,n,r}:=
\operatorname{Led}(x_{\varepsilon,r-1}(x))^\frown
\ell_{\varepsilon,n,r}\). Se construye entonces la tupla
\begin{equation}
 y_{\varepsilon,r}(x):=
 x_{\varepsilon,r-1}(x)
 [\tau\leftarrow M_{\rm ph}(r),\
  h\leftarrow h_r(h(x_{\varepsilon,r-1}(x))),\
  o\leftarrow o(M_{\rm ph}(r)),\
  b_{\rm ref}\leftarrow\mathfrak p_\varepsilon],
 \label{eq:tupla-sel-cal-ch}
\end{equation}
donde la selección deja intactos célula, cocientes, acarreo, ruta, memoria,
cilindro, frontera, supervivencia y registro; la coordenada
\(b_{\rm ref}\) hace explícita la reducción balanceada escogida. No se presupone aquí
pertenencia a una relación ambiental.

A continuación se define \(z_{\varepsilon,r}(x)\) campo por campo: su
configuración es \(\widetilde q_{n,r+1}\); su célula completa es
\(a_{(n,r)^+}\), con los
cuatro residuos y cocientes recalculados por las dos divisiones euclídeas;
conserva la hoja seleccionada y la orientación; y efectúa las seis
actualizaciones siguientes, en las que sólo se omite el argumento común \(x\):
\begin{equation}
\begin{aligned}
 C(z_{\varepsilon,r})&:=C(y_{\varepsilon,r})^\frown
       e_{\varepsilon,n,r},\\
 \partial C(z_{\varepsilon,r})&:=\partial C(y_{\varepsilon,r})^\frown
       (\widetilde q_{n,r},\widetilde q_{n,r+1}),\\
 p_{\APP}(z_{\varepsilon,r})&:=p_{\APP}(y_{\varepsilon,r})^\frown
       a_{(n,r)^+},\\
 p_{\TRIT}(z_{\varepsilon,r})&:=p_{\TRIT}(y_{\varepsilon,r})^\frown
       \tau(r),\\
 \operatorname{Led}(z_{\varepsilon,r})&:=
       \operatorname{Led}^{\rm new}_{\varepsilon,n,r},\\
 \lambda(z_{\varepsilon,r})&:=e_{\varepsilon,n,r}
       \lambda(y_{\varepsilon,r}),
\end{aligned}
 \label{eq:tupla-tra-cal-ch}
\end{equation}
y mantiene todavía los campos
\(c,\operatorname{gr}_{\mathsf M},\operatorname{mem},s\) de
\(y_{\varepsilon,r}(x)\). Su firma, en cambio, se vuelve a tipar en la
configuración de destino mediante la prefirma transportada
\begin{equation}
\begin{aligned}
 \operatorname{Sig}(z_{\varepsilon,r}(x))
  :=\operatorname{Sig}^{\rm cal}_{\widetilde q_{n,r+1}}
       \Bigl(&\widehat h(z_{\varepsilon,r}(x)),
             o(z_{\varepsilon,r}(x)),c(y_{\varepsilon,r}(x)),
             \lambda(z_{\varepsilon,r}(x)),\\[-2pt]
             &\partial C(z_{\varepsilon,r}(x)),
             \jmath_{d_{n,r}}
                (\operatorname{mem}_{d_{n,r}}
                   (y_{\varepsilon,r}(x)))\Bigr).
\end{aligned}
\label{eq:prefirma-transporte-cal-ch}
\end{equation}
Así, y sólo en
\(\mathsf{Tra}^{\rm cal}\), se añaden la arista, sus extremos, la marca
\((\mathfrak p_\varepsilon,r,d_{\varepsilon,r})\) y sus cociclos al registro de
incidencias.

El estado inicial satisface
\[
 c(x_\star)=0=\operatorname{Carr}
   (\operatorname{id}_{\widetilde q_{0,1}})
   =\operatorname{Carr}(\lambda(x_\star)).
\]
En los estados completos \(x_{\varepsilon,r-1}(x)\), y también después de la
selección \(y_{\varepsilon,r}(x)\), se conserva el invariante tipado
\(c(v)=\operatorname{Carr}(\lambda(v))\), pues
\(\mathsf{Sel}^{\rm cal}\) no modifica ninguno de esos dos campos. El
transporte avanza la ruta pero deja pendiente el acarreo; por ello la tupla de
preactualización satisface exactamente
\[
 c(z_{\varepsilon,r}(x))=c(y_{\varepsilon,r}(x))
 =\operatorname{Carr}(\lambda(y_{\varepsilon,r}(x))).
\]
La aplicación \(\operatorname{Upd}^{\rm cal}\) restaura entonces
\(c(x_{\varepsilon,r}(x))=
\operatorname{Carr}(\lambda(x_{\varepsilon,r}(x)))\) mediante la ley de
composición, sin atribuir ese invariante al estado intermedio \(z\).

Finalmente se aplica la función unaria
\(\operatorname{Upd}^{\rm cal}_{e_{\varepsilon,n,r}}
=\mathsf{Mem}^{\rm cal}_{e_{\varepsilon,n,r}}\) a esa tupla. Su
resultado \(x_{\varepsilon,r}(x)\) viene dado por
\begin{equation}
\begin{aligned}
 c\bigl(x_{\varepsilon,r}(x)\bigr)
   &=\operatorname{Carr}\bigl(\lambda(z_{\varepsilon,r}(x))\bigr)\\
   &=H_{\rm Carr}(\lambda(y_{\varepsilon,r}(x)))
       \bigl(\operatorname{Carr}(e_{\varepsilon,n,r})\bigr)
     +Y_{\rm Carr}(e_{\varepsilon,n,r})
       \bigl(c(y_{\varepsilon,r}(x))\bigr),\\
 \lambda\bigl(x_{\varepsilon,r}(x)\bigr)
   &=\lambda(z_{\varepsilon,r}(x)),\\
 \operatorname{gr}_{\mathsf M}
      \bigl(x_{\varepsilon,r}(x)\bigr)
   &=d_{n,r}+1,\\
 \operatorname{mem}_{d_{n,r}+1}
       \bigl(x_{\varepsilon,r}(x)\bigr)
   &=\mathsf C_{\mathsf M}(e_{\varepsilon,n,r})
       \operatorname{mem}_{d_{n,r}}(z_{\varepsilon,r}(x))
     +\kappa_{\mathsf M}(e_{\varepsilon,n,r}),\\
 s\bigl(x_{\varepsilon,r}(x)\bigr)
   &=s(e_{\varepsilon,n,r})\wedge s(z_{\varepsilon,r}(x)),\\
 \operatorname{Sig}\bigl(x_{\varepsilon,r}(x)\bigr)
   &=\operatorname{Sig}^{\rm cal}_{\widetilde q_{n,r+1}}
       \Bigl(\widehat h(z_{\varepsilon,r}(x)),
             o(z_{\varepsilon,r}(x)),
             c(x_{\varepsilon,r}(x)),
             \lambda(z_{\varepsilon,r}(x)),
             \partial C(z_{\varepsilon,r}(x)),\\[-2pt]
   &\hspace{112pt}
             \jmath_{d_{n,r}+1}
              \bigl(\operatorname{mem}_{d_{n,r}+1}
                 (x_{\varepsilon,r}(x))\bigr)\Bigr).
\end{aligned}
\label{eq:upd-unaria-tres-vueltas-ch}
\end{equation}
Todos los campos no enumerados en
\eqref{eq:upd-unaria-tres-vueltas-ch} coinciden con los de
\(z_{\varepsilon,r}(x)\). En particular, la firma se calcula con la hoja,
orientación, frontera y ruta ya fijadas por \(\mathsf{Sel}^{\rm cal}\) y
\(\mathsf{Tra}^{\rm cal}\), y con el acarreo y la memoria calculados en las
líneas anteriores; no se define circularmente a partir del estado final.
Ni \(\varepsilon\), ni \(d_{\varepsilon,r}\), ni \(u\) son argumentos
adicionales de \(\operatorname{Upd}^{\rm cal}\): la función lee los tres datos de la
arista ya construida.

\paragraph{Inducción simultánea de los niveles.}
Completadas las nueve etapas anteriores para un \(G_n\) ya construido,
se define inmediatamente
\begin{equation}
 \gamma_{\varepsilon,k_n}(x):=x_{\varepsilon,9}(x),\qquad
 G_{n+1}:=\{\gamma_{\varepsilon,k_n}(x):
       x\in G_n,\ \varepsilon\in\TRIT\}.
 \label{eq:injerto-total-tres-vueltas-ch}
\end{equation}
\label{eq:tres-vueltas-tpk-ch}
La base \(G_0=\{x_\star\}\), las tuplas intermedias y
\eqref{eq:injerto-total-tres-vueltas-ch} constituyen una única definición por
recursión sobre \(n\). Por tanto, ningún portador posterior se utiliza para
suponer la existencia de \(G_n\): el nivel siguiente sólo se introduce
después de construir todas sus aristas desde el nivel precedente.

\paragraph{Relación de extensión TPK y pertenencia efectiva.}
Sea \(\mathcal E^{\rm cal}_{\widetilde q}\) la menor colección etiquetada que contiene
\(x_\star\) y, en cada paso de la recurrencia anterior, las cuatro tuplas
\(x_{\varepsilon,r-1}(x),y_{\varepsilon,r}(x),
z_{\varepsilon,r}(x),x_{\varepsilon,r}(x)\) en sus configuraciones de origen
y destino. Sea asimismo
\(\mathcal A^{\rm cal}_{\widetilde q_{n,r}}:=\{a_{n,r}\}\), y póngase
\[
 \widetilde{\mathcal Q}_{\rm obs}:=
   \{\widetilde q_{n,r}:n\ge0,\ 1\le r\le9\},
 \qquad
 \mathcal E^{\rm cal}:=\bigsqcup_{\widetilde q}
     \mathcal E^{\rm cal}_{\widetilde q},
 \qquad
 \mathcal F^{\rm cal}_{\widetilde q}:=
     \pi_{\mathcal F}(\mathcal E^{\rm cal}_{\widetilde q}).
\]
Para cada profundidad \(d\), sea
\(\mathcal E^{\rm cal}_d:=
\operatorname{gr}_{\mathsf M}^{-1}(d)\subset\mathcal E^{\rm cal}\). La
coordenada explícita de grado, y no la mera pertenencia de un representante a
alguna imagen del sistema directo, hace disjuntos estos portadores. Sobre ellos
quedan tipadas las proyecciones
\[
\begin{aligned}
 \operatorname{cfg}&:\mathcal E^{\rm cal}\longrightarrow
     \widetilde{\mathcal Q}_{\rm obs},\\
 g(x)&:=\operatorname{fase}\!
       \bigl(\pi_{\rm obs}(\operatorname{cfg}(x))\bigr)\in C_9,\\
 c&:\mathcal E^{\rm cal}\longrightarrow K_{\rm Carr},\\
 s&:\mathcal E^{\rm cal}\longrightarrow\{\bot,\top\},\\
 \lambda&:\mathcal E^{\rm cal}_{\widetilde q}
       \longrightarrow\mathsf{Hist}^{\rm cal}_{\widetilde q},\\
 \operatorname{gr}_{\mathsf M}&:\mathcal E^{\rm cal}
       \longrightarrow\mathbb N_{\ge1},\\
 \operatorname{mem}_d&:\mathcal E^{\rm cal}_d
       \longrightarrow\mathsf M_d,\\
 \operatorname{mem}(x)&:=
       \jmath_d\bigl(\operatorname{mem}_d(x)\bigr)\in\mathsf M
       \quad(x\text{ de profundidad }d).
\end{aligned}
\]
Aquí \(\operatorname{cfg}(x)\) es la segunda coordenada de la tupla y
\(\operatorname{fase}\) lee la fase del calendario observable. Para las tres
clases de portador se fijan las diagonales
\[
\begin{aligned}
 \Delta_{\APP,\widetilde q}
   &:=\{(a,a):a\in\mathcal A^{\rm cal}_{\widetilde q}\},\\
 \Delta_{{\rm fib},\widetilde q}
   &:=\{(f,f):f\in\mathcal F^{\rm cal}_{\widetilde q}\},\\
 \Delta_{{\rm enr},\widetilde q}
   &:=\{(x,x):x\in\mathcal E^{\rm cal}_{\widetilde q}\}.
\end{aligned}
\]
En particular,
\[
 a_{0,1}\in\mathcal A^{\rm cal}_{\widetilde q_{0,1}},
 \qquad
 x_\star\in\mathcal E^{\rm cal}_{\widetilde q_{0,1}}.
\]
Sobre estos portadores se
definen por sus grafos completos, y no mediante implicaciones necesarias, las
relaciones
\begin{align}
 \mathsf{Sel}^{\rm cal}_{e_{\varepsilon,n,r}}
   &:=\{(x_{\varepsilon,r-1}(x),y_{\varepsilon,r}(x)):
          x\in G_n\},\\
 \mathsf{Tra}^{\rm cal}_{e_{\varepsilon,n,r}}
   &:=\{(y_{\varepsilon,r}(x),z_{\varepsilon,r}(x)):
          x\in G_n\},\\
 \operatorname{Ext}^{\APP,\rm cal}_{e_{\varepsilon,n,r}}
   &:=\{(a_{n,r},a_{(n,r)^+})\},
 \label{eq:ext-app-calibrada-ch}\\
 \operatorname{Ext}^{\rm fib,cal}_{e_{\varepsilon,n,r}}
   &:=\{(\pi_{\mathcal F}x_{\varepsilon,r-1}(x),
          \pi_{\mathcal F}x_{\varepsilon,r}(x)):x\in G_n\}.
 \label{eq:ext-fibra-calibrada-ch}
\end{align}
La extensión enriquecida sincronizada es
\begin{equation}
 \operatorname{Ext}^{\rm enr,cal}_{e_{\varepsilon,n,r}}
 :=\{(x_{\varepsilon,r-1}(x),x_{\varepsilon,r}(x)):
       x\in G_n\}.
 \label{eq:ext-enriquecida-calibrada-ch}
\end{equation}
La tercera aplicación queda definida por
\begin{equation}
 \operatorname{graph}
 \bigl(\operatorname{Upd}^{\rm cal}_{e_{\varepsilon,n,r}}\bigr)
 :=\{(z_{\varepsilon,r}(x),x_{\varepsilon,r}(x)):
       x\in G_n\},
 \label{eq:grafo-upd-calibrado-ch}
\end{equation}
donde el segundo componente es exactamente el de
\eqref{eq:upd-unaria-tres-vueltas-ch}. Por consiguiente,
\begin{equation}
 \mathcal U^{(1),\rm cal}_{e_{\varepsilon,n,r}}
 :=\operatorname{Upd}^{\rm cal}_{e_{\varepsilon,n,r}}\circ
   \mathsf{Tra}^{\rm cal}_{e_{\varepsilon,n,r}}\circ
   \mathsf{Sel}^{\rm cal}_{e_{\varepsilon,n,r}}
 =\{(x_{\varepsilon,r-1}(x),x_{\varepsilon,r}(x)):
      x\in G_n\}.
 \label{eq:transicion-cal-en-tpk-ch}
\end{equation}
Cada operador está indexado por la arista
\(e_{\varepsilon,n,r}\) producida por la reducción
\(\mathfrak p_\varepsilon\), y no por un trit suministrado desde fuera. La
coordenada \(b_{\rm ref}\) distingue además sus codominios. Por ello
\(\mathsf{Sel}^{\rm cal}_{e}\),
\(\mathsf{Tra}^{\rm cal}\) y \(\operatorname{Upd}^{\rm cal}\) son funciones
sobre sus respectivos dominios, mientras
la unión
\(\bigcup_{\varepsilon\in\TRIT}
\mathcal U^{(1),\rm cal}_{e_{\varepsilon,n,r}}\) conserva deliberadamente
las tres salidas de la trisección.

Para \(\bullet\in\{\APP,{\rm fib},{\rm enr}\}\), las identidades y las
composiciones no quedan implícitas. Se definen
\begin{align}
 \operatorname{Ext}^{\bullet,\rm cal}_{\operatorname{id}_{\widetilde q}}
   &:=\Delta_{\bullet,\widetilde q},\\
 \operatorname{Ext}^{\bullet,\rm cal}_{e_t\cdots e_1}
   &:=\operatorname{Ext}^{\bullet,\rm cal}_{e_t}\circ\cdots\circ
      \operatorname{Ext}^{\bullet,\rm cal}_{e_1}
 \label{eq:ext-palabra-calibrada-ch}
\end{align}
para toda palabra admisible de aristas. Las dos clausuras de caminos APP y
fibra se forman por separado,
\begin{equation}
 \operatorname{Path}(\operatorname{Ext}^{\APP,\rm cal}),\qquad
 \operatorname{Path}(\operatorname{Ext}^{\rm fib,cal}),
 \label{eq:dos-clausuras-ext-ch}
\end{equation}
y \(\mathsf{Tra}^{\rm cal}\) las sincroniza arista por arista.

Finalmente, los niveles de estados completos y sus truncamientos se fijan
por la propia recurrencia:
\begin{align}
 H^{\rm cal}_{d_{n,r}}
   &:=\{x_{\varepsilon,r-1}(x):
          x\in G_n,\ \varepsilon\in\TRIT\},\\
 H^{\rm cal}_{d_{n,r}+1}
   &:=\{x_{\varepsilon,r}(x):
          x\in G_n,\ \varepsilon\in\TRIT\},\\
 \rho^{\rm cal}_{d_{n,r}+1,d_{n,r}}
   \bigl(x_{\varepsilon,r}(x)\bigr)
   &:=x_{\varepsilon,r-1}(x).
 \label{eq:truncamiento-calibrado-definido-ch}
\end{align}
El último registro conserva \((n,r,\mathfrak p_\varepsilon)\), de modo que
el predecesor del miembro izquierdo es único y \(\rho^{\rm cal}\) está bien
definida. Para \(b>a\), se pone
\(\rho^{\rm cal}_{b,a}:=\rho^{\rm cal}_{a+1,a}\circ\cdots\circ
\rho^{\rm cal}_{b,b-1}\). Esta aplicación recupera el estado anterior
completo, incluida su firma; no pretende invertir una coordenada sobrescrita
sin consultar el registro.

\begin{proposition}[Realización autónoma del TPK mediante transiciones explícitas]
\label{prop:modelo-tpk-calibrado-ch}
La colección
\[
\begin{aligned}
 \mathfrak T_{\rm TPK}^{\rm cal}:=\bigl(
 &(\widetilde q_{n,r})_{n,r},\pi_{\rm obs},
   \widetilde F_{\rm obs},M_{\rm ph},\\
 &(\mathcal A_{\widetilde q}^{\rm cal})_{\widetilde q},
   (\mathcal E_{\widetilde q}^{\rm cal})_{\widetilde q},
   (\mathcal F_{\widetilde q}^{\rm cal})_{\widetilde q},
   \operatorname{Ext}^{\APP,\rm cal},
   \operatorname{Ext}^{\rm fib,cal},
   \operatorname{Ext}^{\rm enr,cal},\rho^{\rm cal},\\
 &\mathsf{Sel}^{\rm cal},\mathsf{Tra}^{\rm cal},
   \operatorname{Upd}^{\rm cal},\mathsf C_{\mathsf M},
   \kappa_{\mathsf M},\operatorname{gr}_{\mathsf M},
   H_{\rm Carr},Y_{\rm Carr},\\
 &\operatorname{Sig}^{\rm cal},\pi_{\mathcal F}\bigr).
\end{aligned}
\]
es un modelo enriquecido del TPK. En particular, cada par del miembro derecho
de \eqref{eq:transicion-cal-en-tpk-ch} es una arista TPK efectiva, y no la
consecuencia inferida de una condición solamente necesaria.
\end{proposition}

\begin{proof}
Las etiquetas \((n,r)\) hacen disjuntas las copias de los portadores y fijan
sin ambigüedad cada origen y destino. La relación
\(\operatorname{Ext}^{\APP,\rm cal}\) transporta la célula completa y
preserva las dos identidades euclídeas; la relación
\(\operatorname{Ext}^{\rm fib,cal}\) relaciona las fibras de los estados
completos anterior y posterior: transporta hoja, orientación, prefijo,
cilindro, frontera y registro e incorpora el acarreo, la memoria y la firma
actualizados. La selección no modifica el registro y el transporte concatena
exactamente una entrada y produce la prefirma de destino
\eqref{eq:prefirma-transporte-cal-ch}. La actualización recalcula la firma mediante la
aplicación tipada \(\operatorname{Sig}^{\rm cal}_{\widetilde q}\) desde
hoja, orientación, acarreo, ruta, frontera y memoria ya determinados. La aplicación
\(\operatorname{Upd}^{\rm cal}\) es unaria porque la arista contenida en
\(z_{\varepsilon,r}(x)\) determina sus cociclos.

Las identidades son las diagonales y las extensiones por palabras son las
composiciones de \eqref{eq:ext-palabra-calibrada-ch}. La concatenación de
prefijos y registros es asociativa; la memoria satisface la ley de cociclo y
el acarreo satisface
\eqref{eq:ley-carry-correcta-ch}. Por tanto, las dos parentizaciones de tres
aristas producen el mismo estado. La fórmula
\eqref{eq:truncamiento-calibrado-definido-ch} define el truncamiento sobre
los estados completos y la unicidad del último registro prueba que recupera
el origen en todos los campos. Éstas son las leyes de identidad, fuente,
destino, composición, memoria, acarreo y truncamiento del TPK enriquecido.
En particular, las relaciones de extensión quedan determinadas como grafos
de aplicaciones efectivas y no como consecuencias de meras condiciones
necesarias.
\end{proof}

Cada \(\gamma_{\varepsilon,k_n}\) comprende exactamente nueve aplicaciones
de \(\operatorname{Upd}^{\rm cal}_{e_{\varepsilon,n,r}}\circ
\mathsf{Tra}^{\rm cal}_{e_{\varepsilon,n,r}}\circ
\mathsf{Sel}^{\rm cal}_{e_{\varepsilon,n,r}}\). La
supervivencia se probará mediante una cola construida y no se utilizará para
definir \(G_n\).

\begin{lemma}[Trisección interna efectiva de una vuelta]
\label{lem:elevacion-catalogal-nonadica}
Para todo \(n\ge0\), todo \(x\in G_n\) y todo
\(\varepsilon\in\{-1,0,+1\}\), la expresión
\(\gamma_{\varepsilon,k_n}(x)\) está definida por nueve transiciones TPK,
con sus relaciones sincronizadas
\(\operatorname{Ext}^{\APP,\rm cal}\) y
\(\operatorname{Ext}^{\rm fib,cal}\), y satisface
\begin{align}
 \rho^{\rm cal}_{k_n+9,k_n}\gamma_{\varepsilon,k_n}(x)&=x,
 \label{eq:truncamiento-injerto-ch}\\
 g\bigl(\gamma_{\varepsilon,k_n}(x)\bigr)
   &=g(x),&
 \jmath_{k_n+9}\!\left(\operatorname{mem}_{k_n+9}
        (\gamma_{\varepsilon,k_n}(x))\right)
   &=\jmath_{k_n}\!\left(\operatorname{mem}_{k_n}(x)\right)+u,
 \label{eq:fase-memoria-injerto-ch}\\
 c\bigl(\gamma_{\varepsilon,k_n}(x)\bigr)
   &=c(x)+\varepsilon\kappa_{\rm Carr}.&&
 \label{eq:carry-injerto-ch}
\end{align}
\label{eq:tres-vueltas-propiedades-ch}
\label{eq:tres-vueltas-acarreo-ch}
El registro ordenado de las nueve aristas determina \(\varepsilon\) de manera única.
Las tres imágenes son, por tanto, disjuntas. Cada extremo pertenece a
\(G_{n+1}\) y vuelve a admitir las tres aplicaciones
\(\gamma_{-1,k_{n+1}},\gamma_{0,k_{n+1}},
\gamma_{+1,k_{n+1}}\).
\end{lemma}

\begin{proof}
Las fórmulas \eqref{eq:celula-completa-fase-ch} recorren una vez las nueve
fases y devuelven la copia \(a_{n+1,1}\) de \(a(1)\). En cada paso, la relación APP transporta el
séxtuplo completo; \(\mathsf{Sel}^{\rm cal}\) fija el tipo TRIT desde la marca de fase
\(r\) conservada en el estado;
y \(\mathsf{Tra}^{\rm cal}\) añade exactamente un símbolo al prefijo, un tramo al
cilindro y su par ordenado de extremos a la frontera. Por inducción sobre
\(r\), las tuplas explícitas \(y_{\varepsilon,r}(x)\) y
\(z_{\varepsilon,r}(x)\) satisfacen los predicados de las relaciones
indicadas, y la actualización unaria produce el estado del nivel siguiente.
La identidad
\(\rho^{\rm cal}_{d_{n,r}+1,d_{n,r}}
(x_{\varepsilon,r}(x))=x_{\varepsilon,r-1}(x)\) es la definición
\eqref{eq:truncamiento-calibrado-definido-ch}; su buena definición procede de
la unicidad del último registro. La composición de las nueve aplicaciones de
truncamiento recupera \(x\), lo que prueba
\eqref{eq:truncamiento-injerto-ch} sin invertir por separado los campos
sobrescritos.

La fase \(9\to1\) aparece una sola vez. Por
\eqref{eq:memoria-arista-ch} y la ley de cociclo, la parte lineal de memoria
transporta la memoria previa y su parte traslacional suma exactamente \(u\)
en el límite directo. La compatibilidad
\(\jmath_{d+1}\jmath_{d+1,d}=\jmath_d\) prueba
\eqref{eq:fase-memoria-injerto-ch}; el retorno concierne a la fase, no al
estado enriquecido.

Por \eqref{eq:carry-arista-derivado-ch}, ocho aristas poseen incremento nulo y
la arista que completa \(\mathfrak p_\varepsilon\) posee el incremento calculado
\(\chi(\varepsilon,\varepsilon)\kappa_{\rm Carr}\). La ley
\eqref{eq:ley-carry-correcta-ch} y
\eqref{eq:carry-no-parametrico-ch} dan
\eqref{eq:carry-injerto-ch}. El registro conserva
\((\mathfrak p_\varepsilon,2\varepsilon,r_3(2\varepsilon),
\chi(\varepsilon,\varepsilon))\); los tres valores de este registro son
distintos incluso cuando el incremento escalar es cero. Por ello las imágenes
no se identifican y el decodificador de bloque recupera una única
\(\varepsilon\).

Las definiciones de las aristas dependen sólo de la nueva fase y concatenan,
en vez de sustituir, prefijo, cilindro, frontera y registro. También son
invariantes bajo
\(\jmath_d\operatorname{mem}_d\mapsto
  \jmath_d\operatorname{mem}_d+nu\)
y bajo la traslación del acarreo de
entrada. En consecuencia, la misma construcción se aplica al extremo de
cualquier vuelta. Para cada estado finito, la sucesión explícita
\[
 x,\quad\gamma_{0,k_n}(x),\quad
 \gamma_{0,k_{n+1}}\gamma_{0,k_n}(x),\quad\ldots
\]
es una cola infinita compatible. Esto demuestra la supervivencia de todos los
estados construidos y la disponibilidad reiterada de las tres ramas, sin
postular ninguna de ambas propiedades. Por consiguiente podemos, finalmente,
identificar, de acuerdo con las definiciones de nivel anteriores,
\begin{equation}
 H_{k_n}^{\rm cal}:=G_n
 \label{eq:subarbol-calibrado-superviviente-ch}
\end{equation}
para todo \(n\ge0\).
\end{proof}

El cierre anterior distingue los tres relojes implicados. Nueve emisiones
devuelven la fase en \(C_9\) y avanzan la memoria. Doce vueltas, es decir,
\(108\) emisiones, devuelven además la posición del calendario observable;
tampoco entonces se borran el registro ni la memoria. El operador
\(\mathcal K_{\rm ph}\) actúa únicamente después, como reconocimiento de
memoria reducida de las vueltas ya construidas; nunca genera las aristas de
\(\operatorname{Ext}^{\rm enr,cal}\).


Sea
\[
 \upsilon:\mathbb F_3\longrightarrow\{-1,0,+1\},\qquad
 \upsilon(0)=0,\quad\upsilon(1)=+1,\quad\upsilon(2)=-1
\]
el representante balanceado. Para
\(b=(b_1,\ldots,b_6)\in\mathbb F_3^6\), la prolongación de bloque es la
composición ordenada de seis vueltas
\[
 \Gamma_{b,k}^{[54]}:=
 \gamma_{\upsilon(b_6),k+45}\circ\gamma_{\upsilon(b_5),k+36}\circ
 \gamma_{\upsilon(b_4),k+27}\circ\gamma_{\upsilon(b_3),k+18}\circ
 \gamma_{\upsilon(b_2),k+9}\circ\gamma_{\upsilon(b_1),k}.
\tag{39d}
\label{eq:gamma-bloque-54-ch}
\]
En el nivel de bloque se emplea siempre
\(k=k_{6N}=1+54N\); por ello cada factor tiene el dominio construido por el
factor anterior.
Fijemos una semilla APP canónica \(x_\star\in H^{\rm cal}_1\) y definamos
recursivamente
\[
 \begin{aligned}
  Y_0&:=\{x_\star\},\qquad Y_N\subseteq H_{1+54N}^{\rm cal},\\
  Y_{N+1}&:=\{\Gamma_{b,1+54N}^{[54]}(x):
                 x\in Y_N,\ b\in\mathbb F_3^6\},\\
  \rho^{\rm blk}_{N+1,N}
    &:=\rho^{\rm cal}_{1+54(N+1),1+54N}
       \!\upharpoonright Y_{N+1},\\
  \operatorname{Ext}^{\rm cal}_N(x)
    &:=\{\Gamma_{b,1+54N}^{[54]}(x):b\in\mathbb F_3^6\}.
 \end{aligned}
\tag{39e}
\label{eq:sincronizacion-54-ch}
\]
La notación \(\rho^{\rm cal}_{b,a}\) designa la composición de todas las supresiones
entre los niveles \(b\) y \(a\). Por construcción, cada prolongación de
\(\operatorname{Ext}^{\rm cal}_N(x)\) es una cadena de cincuenta y cuatro
aristas TPK, no un salto del cociente de Markov.

El lector se define directamente sobre el registro, antes de apelar a una
representación mediante \(\Gamma_b\). El lema precedente proporciona un
decodificador
\[
 \operatorname{dec}_9:\{\text{segmentos generados de nueve aristas}\}
 \longrightarrow\{-1,0,+1\},
\]
que lee la pareja TRIT y su reducción en el segmento. Si
\(\operatorname{Led}^{(9)}_{j,i}(x)\) denota el \(i\)-ésimo segmento nonádico
del bloque \(j\) del registro de \(x\), se fija
\begin{equation}
 \beta_{{\rm blk},N}(x):=
 \Bigl(
   \bigl(\upsilon^{-1}\operatorname{dec}_9
          (\operatorname{Led}^{(9)}_{j,i}(x))\bigr)_{i=1}^{6}
 \Bigr)_{j=0}^{N-1}.
 \label{eq:lector-bloques-ch}
\end{equation}
Para \(N=0\) el valor es la palabra vacía. Esta definición no presupone que
la descomposición de \(x\) como composición de aplicaciones \(\Gamma_b\) sea
única.

Especializamos ahora la conjugación aritmética precedente a
\(p=3\), \(d=6\) y \(A=I_6\). Si
\(\pi_{N+1,N}:(\mathbb Z/3^{N+1}\mathbb Z)^6\to
(\mathbb Z/3^N\mathbb Z)^6\) es la reducción canónica, definimos
\begin{equation}
 \Phi_N:=\Psi_N^{-1}\circ\beta_{{\rm blk},N}:
 Y_N\longrightarrow(\mathbb Z/3^N\mathbb Z)^6.
 \label{eq:phi-hensel-bloques-ch}
\end{equation}
Esta definición no asigna una coordenada después de observar una rama: el
lector \(\beta_{{\rm blk},N}\) se obtiene del registro TPK ya construido y
\(\Psi_N^{-1}\) es la inversa explícita, fijada con anterioridad, de la torre
aritmética no filtrada.

\begin{proposition}[Conjugación finita y cobertura completa de la fibra]
\label{prop:cobertura-729-seccion}
Para todo \(N\geq0\),
\[
 \beta_{{\rm blk},N}:Y_N\xrightarrow{\ \cong\ }
 (\mathbb F_3^6)^N
\]
es una biyección compatible con los truncamientos, y para cada \(x\in Y_N\)
\[
 \begin{gathered}
 \beta_{{\rm blk},N+1}
 \bigl[\operatorname{Ext}^{\rm cal}_N(x)\bigr]
 =\{\beta_{{\rm blk},N}(x)^\frown b:
       b\in\mathbb F_3^6\},\\
 \operatorname{pref}_N\circ\beta_{{\rm blk},N+1}
 =\beta_{{\rm blk},N}\circ\rho^{\rm blk}_{N+1,N}.
 \end{gathered}
\tag{40$_0$}
\label{eq:cobertura-729-ch}
\]
Las aplicaciones \(\Phi_N\) son biyecciones y satisfacen la identidad de
entrelazamiento
\begin{equation}
 \pi_{N+1,N}\circ\Phi_{N+1}
 =\Phi_N\circ\rho^{\rm blk}_{N+1,N}.
 \label{eq:entrelazamiento-hensel-ext-ch}
\end{equation}
Por consiguiente, la torre de historias generadas y la torre aritmética de
Hensel son isomorfas como sistemas inversos finitos.
\end{proposition}

\begin{proof}
La afirmación es inmediata para \(N=0\). Supongamos que vale en el nivel
\(N\). El lema anterior construye, para cada palabra
\(b\in\mathbb F_3^6\), una prolongación
    \(\Gamma_{b,1+54N}^{[54]}(x)\). Si \(b\ne b'\), la primera coordenada en la
que difieren determina una pareja TRIT y un cociclo distintos en el primer
segmento nonádico divergente del registro genealógico. Ese registro conserva
las aristas ordenadas, no sólo la suma final de sus acarreos; por tanto, las dos
historias son distintas incluso si acarreos de vueltas posteriores se
compensan. La lectura directa \eqref{eq:lector-bloques-ch} da entonces
\begin{equation}
 \beta_{{\rm blk},N+1}
 \bigl(\Gamma_{b,1+54N}^{[54]}(x)\bigr)
 =\beta_{{\rm blk},N}(x)^\frown b.
 \label{eq:beta-concatenacion-demostrada-ch}
\end{equation}
Las seis trisecciones sucesivas producen exactamente \(3^6\) prolongaciones
distintas. La supresión de los últimos cincuenta y cuatro pasos recupera
\(x\), y la supresión del último bloque del registro recupera
\(\beta_{{\rm blk},N}(x)\). Esto prueba simultáneamente la biyección y la
compatibilidad de \(\beta_{{\rm blk},N}\) con los truncamientos.

La proposición~\ref{prop:ch-hensel-bruto} prueba, independientemente, que
\(\Psi_N\) es biyectiva y que
\[
 \operatorname{pref}_N\circ\Psi_{N+1}
 =\Psi_N\circ\pi_{N+1,N}.
\]
Al componer esta igualdad con la compatibilidad de \(\beta_{\rm blk}\) y
aplicar \(\Psi_N^{-1}\), se obtiene
\[
 \pi_{N+1,N}\Psi_{N+1}^{-1}\beta_{{\rm blk},N+1}
 =\Psi_N^{-1}\beta_{{\rm blk},N}\rho^{\rm blk}_{N+1,N},
\]
que es exactamente \eqref{eq:entrelazamiento-hensel-ext-ch}. La igualdad no
procede de definir una coordenada ad hoc sobre cada rama: resulta de componer
dos sistemas de biyecciones construidos y compatibles a todas las
profundidades.
\end{proof}

El límite inverso
\[
 \Omega_\Gamma^{\rm cal}:=
 \varprojlim_N(Y_N,\rho^{\rm blk}_{N+1,N})
\tag{40a$_0$}
\label{eq:omega-gamma-catalogal-ch}
\]
es el límite de historias del modelo TPK construido arista por
arista en el lema precedente. Su proyección por \(\pi_{\rm obs}\) conserva
las configuraciones del calendario observable, pero la afirmación no depende
de identificar este modelo autónomo con una relación ambiental definida sólo
por condiciones necesarias. Las biyecciones finitas pasan
al límite y producen el homeomorfismo cilíndrico
\[
 \Lambda:(\mathbb F_3^6)^{\mathbb N}
 \xrightarrow{\ \cong\ }\Omega_\Gamma^{\rm cal},
 \qquad
 \beta_{\rm blk}:=\Lambda^{-1}.
\tag{40b$_0$}
\label{eq:homeomorfismo-calibrado-ch}
\]
La proyección canónica se denota
\(\rho^{\rm blk}_{\infty,N}:\Omega_\Gamma^{\rm cal}\to Y_N\).
Por tanto, \(\Lambda\) establece una equivalencia compatible con los
truncamientos: cada cinta ternaria codifica una historia del TPK y cada
prefijo se prolonga mediante aristas del mismo generador.

Los subcilindros se definen ahora como subconjuntos del límite, no como pares
de sufijos. Para \(x\in Y_N\) y \(b\in\mathbb F_3^6\), sean
\begin{align}
 \mathscr Z_N(x)
   &:=\{\xi\in\Omega_\Gamma^{\rm cal}:
       \rho^{\rm blk}_{\infty,N}(\xi)=x\},
 \label{eq:cilindro-real-limite-ch}\\
 \mathscr Z_{N+1}(x;b)
   &:=\{\xi\in\Omega_\Gamma^{\rm cal}:
       \rho^{\rm blk}_{\infty,N+1}(\xi)
       =\Gamma_{b,1+54N}^{[54]}(x)\}.
 \label{eq:subcilindro-real-limite-ch}
\end{align}

\begin{theorem}[Partición efectiva en \(729\) subcilindros de refinamiento]
\label{thm:particion-729-subcilindros-ch}
Para todo \(N\ge0\) y todo \(x\in Y_N\), se tiene la unión disjunta
\begin{equation}
 \mathscr Z_N(x)=
 \bigsqcup_{b\in\mathbb F_3^6}\mathscr Z_{N+1}(x;b).
 \label{eq:particion-real-729-ch}
\end{equation}
Los \(3^6=729\) términos son subcilindros de refinamiento no vacíos y
clopen; cada uno está realizado por cincuenta y cuatro aristas del TPK
construido.
\end{theorem}

\begin{proof}
Sea \(\xi\in\mathscr Z_N(x)\). Su coordenada de nivel \(N+1\) pertenece a
\(Y_{N+1}\), trunca a \(x\) y, por la biyección
\(\beta_{{\rm blk},N+1}\), tiene un último bloque único
\(b\in\mathbb F_3^6\). La identidad
\eqref{eq:beta-concatenacion-demostrada-ch} implica entonces
\(\rho^{\rm blk}_{\infty,N+1}(\xi)
=\Gamma_{b,1+54N}^{[54]}(x)\); por tanto, \(\xi\) pertenece a uno de los
términos del miembro derecho. La inclusión inversa es inmediata al truncar.
Si dos términos se intersecaran, sus coordenadas de nivel \(N+1\) serían
iguales y el decodificador del registro daría \(b=b'\), de modo que la unión
es disjunta.

Cada extremo \(\Gamma_b^{[54]}(x)\) admite una cola infinita obtenida
iterando la prolongación de bloque nulo; por ello todos los términos son no
vacíos. Como \(Y_{N+1}\) es finito y discreto, son preimágenes de puntos por
una proyección continua del límite inverso y, por tanto, clopen. La propia
definición de \(\Gamma_b^{[54]}\) acredita sus cincuenta y cuatro aristas.
\end{proof}

La colección \((\Phi_N)_N\) induce finalmente el homeomorfismo de límites
\begin{equation}
 \Phi_\infty:\Omega_\Gamma^{\rm cal}
 \xrightarrow{\ \cong\ }\mathbb Z_3^6,\qquad
 \Phi_\infty(\xi):=
 \bigl(\Phi_N(\rho^{\rm blk}_{\infty,N}(\xi))\bigr)_{N\ge0},
 \label{eq:phi-infinito-hensel-ext-ch}
\end{equation}
donde
\(\mathbb Z_3^6=\varprojlim_N(\mathbb Z/3^N\mathbb Z)^6\).
La ecuación \eqref{eq:entrelazamiento-hensel-ext-ch} garantiza que el miembro
derecho es una clase compatible, y las biyecciones finitas prueban
inyectividad, sobreyectividad y continuidad en ambos sentidos. Bajo
\(\Phi_\infty\), \(\mathscr Z_N(x)\) es exactamente la clase de congruencia
\[
 \{z\in\mathbb Z_3^6:z\equiv\Phi_N(x)\pmod{3^N}\},
\]
y los conjuntos \(\mathscr Z_{N+1}(x;b)\) son sus \(729\) clases de refinamiento
módulo \(3^{N+1}\). Éste es el empalme efectivo entre prolongación TPK,
cilindros del límite y torre de Hensel.

Para \(b=(b_1,\ldots,b_6)\in\mathbb F_3^6\), póngase
\[
 \nu(b):=\sum_{i=1}^{6}[b_i]_3\,3^{6-i}
 \in\{0,\ldots,728\},
\tag{40c$_0$}
\label{eq:valor-bloque-ch}
\]
donde \([b_i]_3\in\{0,1,2\}\) es el representante posicional. Las
\(243\) palabras visibles del censo inicial son la imagen de la primera
emisión; no se confunden con la capacidad \(3^6\) de cada bloque de
prolongación. Finalmente,
\[
 X_{\infty,\mathbb Z}^{\rm cal}
 :=\mathbb Z_{\rm HMT}\times\Omega_\Gamma^{\rm cal}.
\]


Para un prefijo \(s=(m;b_0,\ldots,b_{N-1})\), defínanse en la aritmética
generada de \(\mathbb Q_{\rm HMT}\)
\[
 a_N(s):=m+\sum_{q<N}\frac{\nu(b_q)}{729^{q+1}},
 \qquad
 \operatorname{Succ}(s):=
 \{(m;b_0,\ldots,b_{N-1},b):b\in\mathbb F_3^6\}.
\]
Si \((m;\beta_{{\rm blk},N}(x))\) denota el prefijo obtenido anteponiendo
\(m\) a la tupla \(\beta_{{\rm blk},N}(x)\),
\eqref{eq:cobertura-729-ch} da la igualdad tipada
\[
 \operatorname{Succ}(m;\beta_{{\rm blk},N}(x))
 =\{(m;\beta_{{\rm blk},N+1}(y)):
       y\in\operatorname{Ext}^{\rm cal}_N(x)\}.
\]
Los cilindros semiabiertos son, desde su primera definición, subconjuntos de
la completación HMT:
\[
 I_N^{\rm HMT}(s):=
 \{r\in\mathbb R_{\rm HMT}:
      a_N(s)\le_{\rm HMT}r<_{\rm HMT}a_N(s)+729^{-N}\}.
\]
Satisfacen materialmente
\[
 I_N^{\rm HMT}(s)=
 \bigsqcup_{y\in\operatorname{Succ}(s)}I_{N+1}^{\rm HMT}(y),
 \qquad |\operatorname{Succ}(s)|=729,
 \qquad \operatorname{diam}_{\rm HMT}I_N^{\rm HMT}(s)
      =729^{-N}\longrightarrow0.
\tag{40a}\label{eq:ch-particion-cilindros}
\]
La convención de frontera identifica las dos escrituras sólo después de
construirlas. La partición~\eqref{eq:ch-particion-cilindros}, el anidamiento y la contracción producen

\[
\begin{aligned}
 P_{\rm ar}&:X_{\infty,\mathbb Z}^{\rm cal}
     \twoheadrightarrow\mathbb R_{\rm HMT},\\
 P_{\rm ar}(m,\xi)
   &:=\text{el único elemento de }
     \bigcap_N C_N^{\rm cl}
       \bigl(m;\beta_{{\rm blk},N}(\rho^{\rm blk}_{\infty,N}(\xi))\bigr),\\
 C_N^{\rm cl}(s)&:=\overline{I_N^{\rm HMT}(s)}.
\end{aligned}
\tag{40}\label{eq:ch-proyeccion-arquimediana}
\]

Por tanto,

\[
 X_{\infty,\mathbb Z}^{\rm cal}/\!\sim_{P_{\rm ar}}
 \;\cong\;\mathbb R_{\rm HMT}.
\tag{41}\label{eq:ch-cociente-recta}
\]

La otra representación \(P_{\rm exc}\) conserva incidencia, frontera y memoria
hacia Paley--Witt--Golay--Leech y, desde allí, hacia
\(V^\natural\)--Monster. Ambas proceden del
mismo antecedente de~\eqref{eq:ch-limite-inverso}; ni la coordenada arquimediana ni el reconocimiento
excepcional seleccionan las semillas. Las cinco construcciones
\(sp,sol,gau,coh,det\) permanecen consustanciales dentro de
\(\mathcal C_{\rm cont}^{\rm disc}\) y no se disgregan para efectuar el paso
cardinal.

\subsubsection{Descenso de operaciones al cociente de evaluación}
\label{sec:ch-descenso-fibras}

El cociente anterior identifica historias que expresan el mismo valor, no
sus registros completos. El criterio siguiente caracteriza exactamente
cuándo una operación de historias induce una operación en ese cociente sin
escoger una historia privilegiada dentro de cada fibra; su formulación
continúa la construcción genealógica establecida en las secciones anteriores.

\begin{proposition}[Criterio de descenso por fibras de evaluación]
\label{prop:ch-descenso-fibras}
Sea \(q:\Omega\to J\) la evaluación sobre su imagen de una coordenatización
genealógica. Sea \(\star:D\to\Omega\) una operación con
\(D\subseteq\Omega^2\) saturado por las fibras de \(q\times q\).
Existe una única operación
\(\overline\star:(q\times q)(D)\to J\) tal que
\[
 q(x\star y)=q(x)\,\overline\star\,q(y)
\]
si y sólo si, para pares del dominio,
\[
 q(x)=q(x'),\quad q(y)=q(y')
 \quad\Longrightarrow\quad q(x\star y)=q(x'\star y').
\]
\end{proposition}

\begin{proof}
Si existe la operación visible, su resultado depende sólo de los dos
valores, lo que prueba la necesidad. Para la suficiencia, se define
\(s\,\overline\star\,t=q(x\star y)\), donde \(q(x)=s\) y
\(q(y)=t\). La saturación conserva el dominio al cambiar representantes;
la condición enunciada conserva el resultado. La sobreyectividad de
\(q\times q\) sobre la imagen del dominio prueba la unicidad.
\end{proof}

Una coordenada no reemplaza automáticamente una historia como argumento de
todas sus operaciones: el criterio determina cuándo esa sustitución es
legítima. En particular, transportar una operación a través de una biyección
no basta para identificarla con una operación nativa anterior a la lectura.
Si una coordenatización \(\Omega\) es compacta y \(q\) es continua y sobreyectiva
sobre un espacio Hausdorff \(J\), la evaluación induce además un
homeomorfismo \(\Omega/{\sim_q}\to J\): es una biyección continua de un
compacto a un Hausdorff. Es el \emph{cociente}, y no necesariamente el
espacio de historias, el que se identifica topológicamente con \(J\).


\subsubsection{Reconstrucción interna de la coordenatización y levantamiento uniforme}

Dentro de \(\mathfrak G_{\rm HMT}(h,m_\infty)\), la reconstrucción de
\eqref{eq:ch-limite-inverso}--\eqref{eq:ch-cociente-recta} se realiza con los mismos códigos y operadores. Denotemos por
\(\mathbb R_{\rm HMT}^{\mathfrak G}\) la completación HMT interna. La
unicidad del cuerpo ordenado completo arquimediano proporciona, \textbf{después} de
construirlo, el isomorfismo interno
\[
 \eta:\mathbb R_{\rm HMT}^{\mathfrak G}
 \xrightarrow{\ \cong\ }\mathbb R^{\mathfrak G}.
\tag{41$_0$}\label{eq:ch-isomorfismo-interno}
\]
Para el intervalo interno correspondiente
\(I_N^{\mathfrak G}(s):=
 I_N^{\rm HMT}(s)\cap\mathbb R_{\rm HMT}^{\mathfrak G}\), el isomorfismo
preserva extremos y orden:
\[
 \eta[I_N^{\mathfrak G}(s)]
 =\{r\in\mathbb R^{\mathfrak G}:
      a_N(s)\le r<a_N(s)+729^{-N}\}.
\tag{41$_0$a}\label{eq:ch-intervalos-internos}
\]
La sobreyectividad de~\eqref{eq:ch-proyeccion-arquimediana} cubre todos los elementos de esa completación, no
sólo los obtenidos en una ejecución finita.

Para que «todos» no quede oculto en la definición del codominio, se construyen
las dos inyecciones que relacionan la recta con su potencia interna. Sea

\[
 E(A)=\sum_{n\ge0}\frac{2\mathbf1_A(n)}{3^{2n+2}}
 \qquad
 \bigl(A\in\mathcal P^{\mathfrak G}(\omega)\bigr).
\tag{41a}\label{eq:ch-incrustacion-cantor}
\]

La escritura coloca \(0\) o \(2\) únicamente en las posiciones ternarias
pares; su imagen está contenida en \([0,1/4]\). Por tanto nunca alcanza el
extremo derecho del cilindro raíz y dos subconjuntos distintos tienen una
primera posición par distinta: \(E\) es inyectiva y no sufre la ambigüedad de
frontera. Escribamos \(E_{\rm HMT}:=\eta^{-1}\circ E\).

Dado \(A\), se parte del cilindro raíz que contiene \(E_{\rm HMT}(A)\). Si
\(s_N(A)\) es el prefijo ya elegido, escríbase
\(s_N(A)=(0;b_0(A),\ldots,b_{N-1}(A))\). La partición~\eqref{eq:ch-particion-cilindros}, transportada
por \(\eta^{-1}\), entrega un único subcilindro semiabierto
\(s_{N+1}(A)\) que contiene \(E_{\rm HMT}(A)\). Por~\eqref{eq:cobertura-729-ch} esa prolongación es
también una prolongación TPK admisible. Por inducción,
\[
 s_{N+1}(A)\!\upharpoonright N=s_N(A),\qquad
 E_{\rm HMT}(A)\in I_N^{\rm HMT}(s_N(A))\quad(N<\omega).
\tag{41b$_0$}\label{eq:ch-prefijos-uniformes}
\]
La recursión es una fórmula uniforme con parámetros \(A,h\). Como ambos son
internos, Reemplazo en el modelo del lema~\ref{thm:ch-6b} forma la sucesión completa de
bloques \((b_q(A))_{q<\omega}\). Defínase, sin confundir la cinta con la
sucesión de prefijos,
\[
 \xi_A:=\Lambda\bigl((b_q(A))_{q<\omega}\bigr)
 \in\Omega_\Gamma^{\rm cal,\mathfrak G}.
\]
Se obtiene así el \textbf{levantamiento uniforme de la incrustación de Cantor}

\[
 Q:\mathcal P^{\mathfrak G}(\omega)
 \longrightarrow X_{\infty,\mathbb Z}^{\rm cal,\mathfrak G},
 \qquad
 Q(A):=(0,\xi_A),\qquad
 P_{\rm ar}(Q(A))=E_{\rm HMT}(A).
\tag{41b}\label{eq:ch-levantamiento-uniforme}
\]

No se llama a \(Q\) inverso de \(P_{\rm ar}\): su composición expresa la
incrustación \(E_{\rm HMT}\), no la identidad en toda la recta. En la otra
dirección, sea \((q_n)_{n<\omega}\) la enumeración generada de
\(\mathbb Q_{\rm HMT}\) y defínase
\[
 C(r):=\{n<\omega:\eta(q_n)<r\}
 \qquad(r\in\mathbb R^{\mathfrak G}).
\]
El corte racional pertenece a \(\mathcal P^{\mathfrak G}(\omega)\) y
\(C\) es inyectivo por densidad. Las inyecciones \(E\) y \(C\), seguidas de
Cantor--Bernstein interno, prueban

\[
 P_{\rm ar}
 \bigl[X_{\infty,\mathbb Z}^{\rm cal,\mathfrak G}\bigr]
 =\mathbb R_{\rm HMT}^{\mathfrak G},
 \qquad
 \mathbb R_{\rm HMT}^{\mathfrak G}
 \cong\mathbb R^{\mathfrak G}
 \approx\mathcal P^{\mathfrak G}(\omega).
\tag{41c}\label{eq:ch-recta-potencia}
\]

Aquí un real arbitrario se usa únicamente como variable universal para probar
sobreyectividad de un árbol ya construido. No selecciona APP, TRIT, TPK, las
constantes ni las transiciones de ese árbol.

\subsection{Enumeración por códigos genealógicos y condensación HMT}

La regla~\eqref{eq:ch-operador-def} entrega nombres de derivación a partir de

\[
 (\ulcorner\varphi\urcorner,
  \sigma(p_0),\ldots,\sigma(p_{k-1})).
\tag{42}\label{eq:ch-codigo-nombres}
\]

Para cualquier objeto \(a\) de la jerarquía, sean \(\beta(a)\) su primera
etapa y \(\sigma(a)\) su menor árbol de derivación en el orden recursivo del
lema~\ref{thm:ch-6b}. No se exige que \(\sigma(a)\) sea un natural cuando la etapa es no
numerable. La relación

\[
 a\prec_{\mathfrak G}b
 \quad\Longleftrightarrow\quad
 (\beta(a),\sigma(a))<_{\rm lex}
 (\beta(b),\sigma(b))
\tag{42a}\label{eq:ch-orden-genealogico}
\]

es el buen orden genealógico global \(<_{\mathfrak G}\) inducido por el
constructor. Los nombres repetidos se eliminan tomando el menor; ningún real
ni ninguna enumeración de la recta se inserta como dato.

\begin{lemma}[Numerabilidad interna de los niveles anteriores a \(\omega_1^{\mathfrak G}\)]
\label{thm:ch-7}
Para todo
\(\beta<\omega_1^{\mathfrak G}\), el nivel
\(\mathfrak G_\beta(h,m_\infty)\) es numerable dentro de
\(\mathfrak G_{\rm HMT}(h,m_\infty)\), y posee una sobreyección canónica

\[
 e_\beta:\omega\twoheadrightarrow
 \mathfrak G_\beta(h,m_\infty).
\tag{43}\label{eq:ch-enumeracion-niveles}
\]

\end{lemma}
\begin{proof}
El nivel inicial tiene un código numerable. Si \(e_\beta\)
enumera el nivel \(\beta\), los códigos de fórmula y las tuplas finitas de
parámetros de~\eqref{eq:ch-codigo-nombres} forman un conjunto numerable y enumeran el sucesor. Si
\(\lambda<\omega_1^{\mathfrak G}\) es límite, existe internamente alguna
sobreyeción \(s:\omega\twoheadrightarrow\lambda\); se toma la
\(<_{\mathfrak G}\)-menor, orden que ya quedó construido en el lema~\ref{thm:ch-6b}, y se
diagonalizan \(e_{s(n)}\). Finalmente, entre todas las sobreyecciones así
obtenidas se define \(e_\beta\) como la \(<_{\mathfrak G}\)-menor. La
recursión \(\beta\mapsto e_\beta\) es una clase uniformemente definible de
\(\mathfrak G\); cada \(e_\beta\), y cada restricción de la recursión a un
ordinal, pertenece a \(\mathfrak G\). Sólo usa \(h,u_{h,m}\) y el buen orden
global.
\end{proof}

\begin{lemma}[Condensación de los reales internos por debajo de \(\omega_1^{\mathfrak G}\)]
\label{thm:ch-8}
Todo real de
\(\mathfrak G_{\rm HMT}(h,m_\infty)\) aparece en un nivel de índice menor que
\(\omega_1^{\mathfrak G}\).

\end{lemma}
\begin{proof}
Sea \(r\subseteq\omega\) un real de \(\mathfrak G\) y sea
\(B:=\mathfrak G_0(h,m_\infty)\). Esta base es transitiva y numerable
internamente. Con los códigos de naturales y pares finitos ya fijados,
\(\operatorname{Ord}\cap B=\omega+1\). Por inducción,
\[
 \operatorname{Ord}\cap\mathfrak G_\alpha=(\omega+1)+\alpha.
\]
El sucesor añade el ordinal formado por los ordinales del nivel anterior;
no puede añadir uno mayor, pues no sería un subconjunto de ese nivel. Los
límites son uniones. Elegimos \(\Theta\) límite, \(\Theta>\omega^2\),
con \(r\in\mathfrak G_\Theta\). Así
\(\operatorname{Ord}\cap\mathfrak G_\Theta=\Theta\), y todo nivel
\(\mathfrak G_\xi\), con \(\xi<\Theta\), pertenece a
\(\mathfrak G_\Theta\). Formamos internamente la estructura expandida
\[
 \mathcal A_\Theta=
 \langle\mathfrak G_\Theta,\in,u_{h,m},
         H_\Theta,<_{\mathfrak G}\!\upharpoonright\mathfrak G_\Theta\rangle
\]
con la relación de niveles como predicado de su signatura,
\(H_\Theta=\{\langle\xi,x\rangle:\xi<\Theta\ \&\
x\in\mathfrak G_\xi\}\), y sus funciones de Skolem definibles. Sea
\[
 Y=\operatorname{Hull}^{\mathcal A_\Theta}
   (B\cup\{B,r\})\prec\mathcal A_\Theta .
\]
La clausura bajo la colección numerable de funciones de Skolem muestra dentro
de \(\mathfrak G\) que \(Y\) es numerable. Sea
\(c:Y\rightarrow M\) su colapso de Mostowski y
\[
 \bar\Theta:=\operatorname{otp}(Y\cap\Theta)
 <\omega_1^{\mathfrak G}.
\]
Como \(B\subseteq Y\) es transitivo, la inducción de pertenencia muestra
que \(c\) fija cada elemento de \(B\) y también \(B\). En particular,
fija \(\omega,u_{h,m},h,m_\infty\). Puesto que
\(r\subseteq\omega\subseteq Y\), también \(c(r)=r\).

Probamos por inducción sobre \(\xi\in Y\cap\Theta\) que
\[
 c``\bigl(Y\cap\mathfrak G_\xi(h,m_\infty)\bigr)
 =\mathfrak G_{c(\xi)}(h,m_\infty).
\tag{43a}\label{eq:ch-condensacion}
\]
Para \(\xi\in Y\cap\Theta\), el predicado \(H_\Theta\) define el
conjunto \(\mathfrak G_\xi\) desde \(\xi\), de modo que
\(\mathfrak G_\xi\in Y\). Relativizar cada fórmula a ese conjunto da
\(Y\cap\mathfrak G_\xi\prec\mathfrak G_\xi\), en el lenguaje al que
el lema~\ref{thm:ch-6} traduce las definiciones HMT.

La base de~\eqref{eq:ch-condensacion} ya está probada porque el colapso fija \(B\). En un
sucesor, si \(\xi+1\in Y\), también \(\xi\in Y\) y
\(c(\xi+1)=c(\xi)+1\). Para
\(a\in Y\cap\mathfrak G_{\xi+1}\), fijamos una fórmula finita que lo
define sobre \(\mathfrak G_\xi\). La afirmación
\[
 a=\{z\in\mathfrak G_\xi:
                  \varphi^{\mathfrak G_\xi}(z,\bar p)\}
\]
es una fórmula ordinaria sobre \(\mathfrak G_\Theta\), con parámetros
\(a,\mathfrak G_\xi\). La elementalidad proporciona los parámetros
\(\bar p\) dentro de \(Y\cap\mathfrak G_\xi\). La hipótesis inductiva
y la elementalidad de ese nivel muestran que \(c(a)\) es el subconjunto
definido por la misma fórmula y los parámetros colapsados sobre
\(\mathfrak G_{c(\xi)}\). Recíprocamente, dados una fórmula y parámetros
en el nivel colapsado, se levantan los parámetros al nivel original.
El subconjunto allí definido existe, es único y pertenece a \(Y\) por
elementalidad. Su colapso es el subconjunto solicitado. Esto prueba ambas
inclusiones en el sucesor, sin introducir un predicado uniforme de verdad.

Para \(\lambda\in Y\cap\Theta\) límite, la elementalidad proporciona,
para cada \(x\in Y\cap\mathfrak G_\lambda\), un testigo
\(\xi\in Y\cap\lambda\) con \(H_\Theta(\xi,x)\). Por tanto,
\[
 Y\cap\mathfrak G_\lambda
 =\bigcup_{\xi\in Y\cap\lambda}(Y\cap\mathfrak G_\xi),
 \qquad c``(Y\cap\lambda)=c(\lambda).
\]
La continuidad de la jerarquía y la hipótesis inductiva dan~\eqref{eq:ch-condensacion} en el
límite. No se requiere cofinalidad de \(Y\cap\lambda\) en el ordinal
ambiental \(\lambda\): sus índices colapsados recorren \(c(\lambda)\).

Finalmente, cada \(x\in Y\) posee un testigo \(\xi\in Y\cap\Theta\)
con \(H_\Theta(\xi,x)\), y \(\bar\Theta\) es límite. Tomando uniones,
\[
 M=\bigcup_{\eta<\bar\Theta}\mathfrak G_\eta(h,m_\infty)
  =\mathfrak G_{\bar\Theta}(h,m_\infty).
\]
Como \(c(r)=r\in M\), se concluye
\(r\in\mathfrak G_{\bar\Theta}\) con
\(\bar\Theta<\omega_1^{\mathfrak G}\).
\end{proof}

El lema no dice que un prefijo de longitud finita enumere la recta. Dice que
cada objeto hereditario generado posee un primer nivel transfinito y un primer
código genealógico dentro de ese nivel. Ésta es exactamente la información que
el reducto extensional pierde cuando conserva sólo el valor expresado.

\subsection{Inyección de la recta generada en el primer ordinal no numerable de la clausura genealógica}

Sea \(\iota(r)\subseteq\omega\) el corte racional de
\(r\in\mathbb R_{\rm HMT}^{\mathfrak G}\) respecto de la enumeración de
\(\mathbb Q_{\rm HMT}\), construida previamente desde los enteros y cocientes
APP. El mapa \(r\mapsto\iota(r)\) es definible desde \(r\) y esa enumeración,
de modo que \(\iota(r)\in\mathfrak G\); el lema~\ref{thm:ch-8} lo sitúa en una etapa de
índice numerable. Definimos

\[
 \begin{aligned}
  \beta(r)&:=\min\{\beta:
       \iota(r)\in\mathfrak G_{\beta+1}(h,m_\infty)\},\\
  n(r)&:=\min\{n:e_{\beta(r)+1}(n)=\iota(r)\},\\
  j_{\rm HMT}(r)&:=\omega\cdot\beta(r)+n(r).
 \end{aligned}
\tag{44}\label{eq:ch-indices-real}
\]

La definición~\eqref{eq:ch-indices-real} es uniforme en los parámetros \(h,m_\infty\). Como
\(\mathbb R_{\rm HMT}^{\mathfrak G}\) es un conjunto de \(\mathfrak G\),
Separación fija el dominio de la definición y Reemplazo en \(\mathfrak G\)
forma su grafo:
\[
 \{\langle r,\omega\cdot\beta(r)+n(r)\rangle:
 r\in\mathbb R_{\rm HMT}^{\mathfrak G}\}\in\mathfrak G.
\]
Por tanto, \(j_{\rm HMT}\) es una función interna de \(\mathfrak G\).

Por los lemas~\ref{thm:ch-7} y~\ref{thm:ch-8},
\(j_{\rm HMT}(r)<\omega_1^{\mathfrak G}\) para todo
real generado. Si \(j_{\rm HMT}(r)=j_{\rm HMT}(s)\), la división ordinal por
\(\omega\) recupera el mismo par \((\beta,n)\); \eqref{eq:ch-enumeracion-niveles} recupera el mismo corte
racional y, por completitud, \(r=s\). Por ello

\[
 j_{\rm HMT}:\mathbb R_{\rm HMT}^{\mathfrak G}
 \hookrightarrow\omega_1^{\mathfrak G}
\tag{45}\label{eq:ch-inyeccion-j}
\]

es una inyección construida por el rango y el nombre de la genealogía HMT. Su
existencia utiliza exactamente la clausura semántica~\eqref{eq:ch-operador-def}, el buen orden de
nombres del lema~\ref{thm:ch-6b} y la condensación~\eqref{eq:ch-condensacion}; no recibe la hipótesis del continuo para elegir la
etapa. La identificación conjuntista posterior describe esta misma
maquinaria en otra notación, sin invertir su procedencia causal.

\subsection{Inyección del primer ordinal no numerable en la recta y equivalencia cardinal}

Para cada \(\beta<\omega_1^{\mathfrak G}\), el lema~\ref{thm:ch-7} permite codificar en un
subconjunto \(w_\beta\subseteq\omega\) un buen orden numerable de tipo
\(\beta\). Para no confundir el orden vacío con el unipuntual, se codifican
tanto el dominio \(D\subseteq\omega\) como su relación de orden estricto
\(R\subseteq D^2\). Fijado un emparejamiento biyectivo
\(\langle\cdot,\cdot\rangle:\omega^2\to\omega\), el código es
\[
 w(D,R)=\{2n:n\in D\}\cup
        \{2\langle a,b\rangle+1:(a,b)\in R\}.
\]
Los bits pares recuperan el dominio y los impares la relación. Se toma el
\(<_{\mathfrak G}\)-menor real que codifica tal buen orden. Como la propiedad
«codifica un buen orden de tipo \(\beta\)» y el buen orden global son internos,
Reemplazo forma la función \(\beta\mapsto w_\beta\) dentro de
\(\mathfrak G\). Los tipos de orden distintos producen códigos distintos. La
aplicación ternaria, ya usada en~\eqref{eq:ch-incrustacion-cantor}, se define por

\[
 E(w)=\sum_{n\ge0}\frac{2\mathbf1_w(n)}{3^{2n+2}}
\tag{46}\label{eq:ch-incrustacion-ternaria}
\]

y es inyectiva y evita el extremo del cilindro raíz. El levantamiento~\eqref{eq:ch-levantamiento-uniforme}
asegura que \(E_{\rm HMT}(w_\beta)\) pertenece a la recta expresada por una
historia HMT. Así

\[
 k_{\rm HMT}(\beta):=E_{\rm HMT}(w_\beta),
 \qquad
 k_{\rm HMT}:\omega_1^{\mathfrak G}
 \hookrightarrow\mathbb R_{\rm HMT}^{\mathfrak G}.
\tag{47}\label{eq:ch-inyeccion-k}
\]

Las inyecciones~\eqref{eq:ch-inyeccion-j} y~\eqref{eq:ch-inyeccion-k}, construidas sin utilizar la igualdad buscada,
dan por Cantor--Bernstein

\[
 \bigl|\mathbb R_{\rm HMT}^{\mathfrak G}\bigr|
 =\aleph_1^{\mathfrak G}.
\tag{48}\label{eq:ch-igualdad-cardinal}
\]

Como~\eqref{eq:ch-recta-potencia} identifica la recta HMT completa con
\(\mathcal P^{\mathfrak G}(\omega)\), \eqref{eq:ch-igualdad-cardinal} se reescribe exactamente como

\[
 \boxed{
  \mathfrak G_{\rm HMT}(h,m_\infty)
  \models 2^{\aleph_0}=\aleph_1.}
\tag{49}\label{eq:ch-decision}
\]

La igualdad no procede de \(729\), de un producto informal entre cardinalidad
y ordinalidad, de las cifras de una constante ni de una horquilla numérica.
Procede de la cobertura de la recta por el límite inverso, la condensación de
la clausura genealógica y las dos inyecciones explícitas.

\subsection{Cuantificador completo sobre todos los subconjuntos internos}

Sea
\(A\in\mathcal P^{\mathfrak G}
(\mathbb R_{\rm HMT}^{\mathfrak G})\). Si \(A\) es numerable, se satisface la
primera alternativa. Si no lo es, \(j_{\rm HMT}[A]\) es un subconjunto no
numerable de \(\omega_1^{\mathfrak G}\). No puede estar acotado por un ordinal
numerable, pues cada segmento inicial de \(\omega_1\) es numerable. Su
enumeración creciente tiene tipo \(\omega_1^{\mathfrak G}\), y la inversa de
\(j_{\rm HMT}\) sobre su imagen produce una inyección
\(\omega_1^{\mathfrak G}\hookrightarrow A\). La inclusión
\(A\subseteq\mathbb R_{\rm HMT}^{\mathfrak G}\), \eqref{eq:ch-igualdad-cardinal} y
Cantor--Bernstein dan

\[
 \boxed{
 \forall A\in\mathcal P^{\mathfrak G}
   (\mathbb R_{\rm HMT}^{\mathfrak G}),\qquad
 |A|\le\aleph_0
 \quad\text{o}\quad
 |A|=\bigl|\mathbb R_{\rm HMT}^{\mathfrak G}\bigr|.}
\tag{50}\label{eq:ch-dicotomia}
\]

Aquí se cuantifica la potencia completa del continuo integral generado por
HMT. El transporte biyectivo del plegado con memoria lleva~\eqref{eq:ch-dicotomia} a cualquiera
de sus representaciones sin perder partes; no sustituye las inyecciones
\eqref{eq:ch-inyeccion-j}--\eqref{eq:ch-inyeccion-k}.

\subsection{Representación de la clausura genealógica como \texorpdfstring{\(L[h,m_\infty]\)}{L[h,m-infinito]}}

Sólo después de~\eqref{eq:ch-decision}--\eqref{eq:ch-dicotomia} se asigna
a~\eqref{eq:ch-clausura-transfinita} su nombre en la notación
conjuntista usual. Sea \(L[u_{h,m}]\) la jerarquía constructible relativa al
real ya producido en~\eqref{eq:ch-codigo-conjunto}. Como la signatura HMT es numerable y todos sus
operadores son relaciones de grafo codificadas por \(h\), el lema~\ref{thm:ch-6} elimina
sus símbolos uniformemente. Se obtiene la igualdad de clases

% Se separa el contador hipervinculado de la etiqueta editorial (51) para
% evitar dos destinos PDF homónimos en la última ecuación numerada del capítulo.
\stepcounter{equation}
\[
 \mathfrak G_{\rm HMT}(h,m_\infty)
 =L[u_{h,m}]=L[h,m_\infty].
\tag{51}\label{eq:ch-representacion-L}
\]

Para probar la primera inclusión se fortalece la inducción: para cada
ordinal \(\beta\), tanto \(\mathfrak G_\beta\) como la secuencia
\(\langle\mathfrak G_\xi:\xi\leq\beta\rangle\) pertenecen a
\(L[u_{h,m}]\). La base es la clausura transitiva de un conjunto construido
desde \(\omega\) y \(u_{h,m}\); los parámetros \(h,m_\infty\) se recuperan
de \(u_{h,m}\). En el sucesor, el nivel precedente es un \emph{conjunto
interno}, no sólo una subclase de \(L[u_{h,m}]\). Su relación de satisfacción
y la eliminación uniforme~\eqref{eq:ch-eliminacion-uniforme} permiten formar internamente el conjunto de
sus subconjuntos definibles. En un límite, la recursión transfinita,
Reemplazo y Unión de \(L[u_{h,m}]\) forman la secuencia anterior y su unión.
La transitividad da entonces \(\mathfrak G\subseteq L[u_{h,m}]\).

Recíprocamente, el lema~\ref{thm:ch-6b} proporciona un modelo interno
transitivo de ZF, con todos los ordinales y con \(u_{h,m}\) como elemento.
La recursión que construye \(L[u_{h,m}]\) se efectúa en ese modelo: en cada
etapa, satisfacción sobre el nivel anterior y definibilidad relativa son
absolutas entre los dos dominios transitivos. Por inducción, sus niveles
coinciden con los de la construcción ambiental y pertenecen a
\(\mathfrak G\). Así \(L[u_{h,m}]\subseteq\mathfrak G\). Finalmente,
\eqref{eq:ch-codigo-conjunto} hace a \(u_{h,m}\) y al par \((h,m_\infty)\)
primitivo-recursivamente interdefinibles, de lo cual se sigue la segunda
igualdad.

La ecuación~\eqref{eq:ch-representacion-L} identifica explícitamente la
clausura con la jerarquía constructible relativa al parámetro generado. Esta
representación no modifica la procedencia causal: \(L[u]\) ni selecciona el
registro genealógico ni define las prolongaciones o la recta; la decisión de
la hipótesis del continuo resulta de la regla de exhaustividad y de los mapas probados entre
\eqref{eq:ch-clausura-transfinita} y~\eqref{eq:ch-dicotomia}. La composición
\eqref{eq:ch-codigo-conjunto}, \eqref{eq:ch-clausura-transfinita},
\eqref{eq:ch-equivalencia-semantica} y
\eqref{eq:ch-decision}--\eqref{eq:ch-representacion-L}, junto con el
teorema~\ref{thm:decision-continuo-hmt}, ofrece una expresión completa y
uniforme del universo realizado, la coincidencia semántica y el cierre
cardinal; ninguna de estas representaciones sustituye la generación
APP--TRIT--TPK.

\subsection{Teorema de resolución afirmativa de la hipótesis del continuo}

\begin{theorem}[Decisión del continuo en HMT]
\label{thm:decision-continuo-hmt}
Para toda historia realizada
\(m_\infty\) producida por la acción conjunta
\(\mathrm{APP}\to\mathrm{TRIT}\to\mathrm{TPK}\to
X_\infty^{\rm enr}\to\mathcal C_{\rm cont}^{\rm disc}\), la extensión
genealógica HMT~\eqref{eq:ch-clausura-transfinita} genera la reconstrucción interna de la recta integral
\eqref{eq:ch-proyeccion-arquimediana}, \(\mathbb R_{\rm HMT}^{\mathfrak G}\), construye las inyecciones
\eqref{eq:ch-inyeccion-j} y~\eqref{eq:ch-inyeccion-k}, y satisface
\(2^{\aleph_0}=\aleph_1\) sobre toda su potencia interna. Por tanto, HMT
demuestra afirmativamente la hipótesis del continuo para el continuo
íntegro que define y genera. La notación \(L[h,m_\infty]\) es una
representación posterior de esa jerarquía, no la premisa que decide el
resultado.
\end{theorem}


\subsection{Criterios de falsación del cierre cardinal}

La lista siguiente identifica los fallos que invalidarían un paso del cierre
anterior. Se distinguen los controles de construcción, semántica y
cardinalidad:

\begin{enumerate}
\item existe un nivel \(N\) o un estado superviviente sin sucesor compatible en
\eqref{eq:ch-extension-no-vacia};
\item \(u_{h,m}\) no codifica hereditariamente \(h,m_\infty\), o alguna relación
de grafo de \(\Sigma_{\rm HMT}\) no admite la eliminación uniforme~\eqref{eq:ch-eliminacion-uniforme};
\item un objeto admitido por las reglas semánticas HMT carece de término en
\eqref{eq:ch-terminos}, o un término de~\eqref{eq:ch-terminos} expresa un objeto ajeno a esas reglas, rompiendo
\eqref{eq:ch-equivalencia-semantica};
\item \(\mathfrak G\) no satisface alguno de los axiomas usados, falla su buen
orden definible o una supuesta enumeración \(e_\beta\) no es interna y
uniforme;
\item los subcilindros de~\eqref{eq:ch-particion-cilindros} no forman una partición disjunta exhaustiva, sus
cilindros dejan de contraerse o la recursión~\eqref{eq:ch-prefijos-uniformes} no produce una
historia compatible interna;
\item el colapso del lema~\ref{thm:ch-8} no conmuta con una etapa sucesora o límite, o existe
un real interno que no aparece en ningún nivel de índice menor que
\(\omega_1^{\mathfrak G}\);
\item la selección \(\beta\mapsto w_\beta\) no es interna, uniforme o conserva
el tipo de orden;
\item dos reales diferentes tienen el mismo par genealógico \((\beta,n)\) en
\eqref{eq:ch-indices-real}, dos ordinales diferentes producen el mismo código en~\eqref{eq:ch-inyeccion-k}, o falla
alguna de las inyecciones \(E,C,j_{\rm HMT},k_{\rm HMT}\);
\item existe una parte interna no numerable de la recta cuyo rango bajo
\(j_{\rm HMT}\) sea acotado en \(\omega_1^{\mathfrak G}\);
\item alguna etapa generativa anterior a la cobertura~\eqref{eq:ch-particion-cilindros} utiliza la hipótesis del continuo,
\(\mathbb R\), una constante objetivo, el Comité de Datos del Consejo
Internacional de Ciencia (CODATA), el Particle Data Group (PDG) o una identificación física
posterior como antecedente causal para definir una semilla, una relación
\(\mathscr L_\tau\), \(\operatorname{Ext}\), un coeficiente, la admisibilidad
o una prolongación del árbol. No constituye tal inversión el levantamiento
posterior~\eqref{eq:ch-prefijos-uniformes}: allí un real arbitrario ya construido es una variable
universal y selecciona, dentro de la partición exhaustiva previamente
demostrada, el único subcilindro que lo contiene.
\end{enumerate}

Ninguna ejecución limitada a un número prefijado de niveles sustituye el
cuantificador universal del sistema inverso; lo falsificaría una ruptura
efectiva de compatibilidad en alguna profundidad. Ninguna notación conjuntista
posterior prueba por sí sola~\eqref{eq:ch-decision}. El cierre utiliza conjuntamente el
lema~\ref{thm:ch-6b}, la identificación recta--potencia de~\eqref{eq:ch-recta-potencia}, los
lemas~\ref{thm:ch-7}--\ref{thm:ch-8} y las dos inyecciones construidas desde
la genealogía.
